\documentclass[a4paper,fleqn]{cas-sc}

\usepackage{float}
\usepackage[utf8]{inputenc}
\usepackage[numbers]{natbib}
\usepackage{tabularx}
\usepackage{array}
\usepackage[english]{babel}
\usepackage{amsmath}
\usepackage{amsfonts}
\usepackage{amssymb}
\usepackage{booktabs}
\usepackage{algorithm}
\usepackage{algpseudocode}
% Float placement: the manuscript carries 24 tables and 10 figures. Allow pages
% to hold more than two floats and keep text pages from breaking early.
\setcounter{topnumber}{3}
\setcounter{bottomnumber}{2}
\setcounter{totalnumber}{5}
\renewcommand{\topfraction}{0.95}
\renewcommand{\bottomfraction}{0.95}
\renewcommand{\textfraction}{0.05}
\renewcommand{\floatpagefraction}{0.7}
% Tighter float and caption spacing, set at begin-document so the class
% defaults cannot override them.
\AtBeginDocument{%
  \setlength{\textfloatsep}{10pt plus 2pt minus 2pt}%
  \setlength{\floatsep}{8pt plus 2pt minus 2pt}%
  \setlength{\intextsep}{8pt plus 2pt minus 2pt}%
  \setlength{\abovecaptionskip}{6pt}%
  \setlength{\belowcaptionskip}{2pt}%
}
\usepackage{multirow}
\usepackage{hyperref}
\usepackage{color}
\usepackage{xcolor}
\usepackage{orcidlink}
\usepackage{lineno}
\usepackage{silence}
\WarningFilter{hyperref}{Ignoring empty anchor}
\setlength{\emergencystretch}{2em}
\setlength{\textfloatsep}{4pt plus 2pt minus 2pt}
\setlength{\floatsep}{4pt plus 2pt minus 2pt}
\setlength{\intextsep}{4pt plus 2pt minus 2pt}
\setlength{\bibsep}{0.2pt plus 0.1pt minus 0.1pt}
% Table styling: plain grid tables with roomier cells and rows.
\renewcommand{\arraystretch}{1.08}
\setlength{\tabcolsep}{5pt}

% Manuscript-specific markers. Keep them visible: an unmarked placeholder is
% a fabrication risk. \pending{...} is the only permitted placeholder for a
% confirmatory quantity.
\newcommand{\pending}[1]{\textcolor{red}{[PENDING: #1]}}
\newcommand{\tbd}{\textcolor{red}{[TBD]}}

% Figure helper: include the file when it exists, otherwise draw a labelled
% placeholder box. The ten figures are generated by scripts/paper_figures/make_figures.py
% and scripts/paper_figures/make_result_figures.py; the manuscript must stay readable
% (and compilable) before they land.
\newcommand{\figfile}[2][\linewidth]{%
  \IfFileExists{#2}%
    {\includegraphics[width=#1]{#2}}%
    {\fbox{\begin{minipage}[c][5.4cm][c]{0.90\linewidth}%
      \centering\ttfamily\small #2\\[4pt]%
      \normalfont\itshape figure pending\end{minipage}}}%
}

% Preprint footer: drop the "submitted to Elsevier" wording.
\ExplSyntaxOn
\cs_set:Npn \__first_footerline:
  {
    \group_begin:
    \small
    \sffamily
    \ifnum\theblind>0\relax
    \else
    \__short_authors: :~
    \fi
    \rmfamily \itshape Preprint
    \group_end:
  }
\ExplSyntaxOff
% The CAS class draws email/URL icons from a thumbnails/ folder. Uploads that
% do not carry that folder fail with "File `thumbnails/cas-email.jpeg' not
% found". The class has a no-logo mode, so the email footnote prints as plain
% text and the document needs no image assets.
\ExplSyntaxOn
\bool_gset_true:N \g_stm_nologo_bool
\ExplSyntaxOff
% --- ORCID badges in the byline -------------------------------------------
% The CAS class prints ORCID iDs only in the footnote. Below is the class's
% own \author definition (cas-common.sty v2.3) with one addition: when an
% author declares [orcid=...], the green iD badge follows the name before the
% affiliation mark. Every other line is unchanged from the class.
\makeatletter
\ExplSyntaxOn
\RenewDocumentCommand \author { O{} m O{} }
   {
     \ResetMarks
     \tl_if_blank:nTF { #3 } { }
       { \keys_set:nn { stm / author } { #3 } }
%
     \int_gincr:N \g_stm_au_int
     \tex_gdef:D \theau@ { \int_use:N \g_stm_au_int }
%
     \seq_gput_right:Nn \g_stm_prelimsau_seq { #2 }
     \bool_if:NTF \l_stm_au_collab_bool
       { \seq_gput_right:cn { g_stm_clau\int_use:N \g_stm_augr_int _seq } }
       { \seq_gput_right:cn { g_stm_au\int_use:N \g_stm_augr_int _seq } }
       {
         \int_gincr:N \g_stm_aau_int
         \tex_gdef:D \theauthor {\int_use:N \g_stm_aau_int }
         \keys_set:nn { stm /author } { #3 }
         \tl_if_head_eq_catcode:nNTF { #1 } a
            { \processAffRef { #1 } }
            { \processAffNum { #1 } }
         \tl_if_empty:NF \l_stm_au_prefix
            { \l_stm_au_prefix_tl \c_space_token }
         \str_if_eq:VnTF \l_stm_au_style_tl  { chinese }
           {
             \invparsename { #2 }
             \textcolor{\l_stm_augroup_color_tl}{\surname}
             \no_break_space:
             \textcolor{\l_stm_augroup_color_tl !50}{\firstname}
           }
           {
             \parsename { #2 }
             \textcolor{\l_stm_augroup_color_tl !50}{\firstname}
             \no_break_space:
             \textcolor{\l_stm_augroup_color_tl}{\surname}
           }
       \tl_if_empty:NF \l_stm_au_suffix_tl
       { \c_space_token \l_stm_au_suffix_tl }
        \tl_if_empty:NF \l_stm_au_orcid_tl
          { \,\small \exp_args:NV \orcidlink \l_stm_au_orcid_tl }
        \unskip
        \textsuperscript
        {
          \tl_if_blank:nTF { #1 }
          { \tex_def:D \sep{} }
          { {\itshape\listAff} \tex_def:D \sep{\unskip,} }
          \process@marks
          \bool_if:NT \l_stm_au_deceased_bool
          { \sep \maltese
            \tex_def:D \sep { \unksip, }
          }
        }
       \tl_if_empty:NF \l_stm_au_degree_tl
       { ,\c_space_token \l_stm_au_degree_tl }
       \tl_if_empty:NF \l_stm_au_role_tl
       { \c_space_token (\l_stm_au_role_tl) }
%
      \ResetMarks
     }
%
     \bool_if:NT \l_stm_au_deceased_bool
      {
       \seq_gput_right:Nn \g_stm_maltese_seq
        {
         \tex_def:D \thefootnote { \maltese }
         \footnotetext{Deceased~author.}
        }
      }
% various social media
     \tl_if_empty:NTF \l_stm_au_facebook_tl { }
       {
         \wrAux { \token_to_str:N \facebookauthor
           { \l_stm_au_facebook_tl } { \exp_not:n {#2} } }
       }
     \tl_if_empty:NF \l_stm_au_twitter_tl
       {
         \wrAux { \token_to_str:N \twitterauthor
           { \l_stm_au_twitter_tl } { \exp_not:n {#2} } }
       }
     \tl_if_empty:NF \l_stm_au_gplus_tl
       {
         \wrAux { \token_to_str:N \gplusauthor
           { \l_stm_au_gplus_tl } { \exp_not:n {#2} } }
       }
     \tl_if_empty:NF \l_stm_au_linkedin_tl
       {
         \wrAux { \token_to_str:N \linkedinauthor
           { \l_stm_au_linkedin_tl } { \exp_not:n {#2} } }
       }
     \tl_if_empty:NF \l_stm_au_orcid_tl
       {
         \wrAux { \token_to_str:N \orcidauthor
           { \l_stm_au_orcid_tl } { \exp_not:n {#2} } }
       }
%
       \@eadauthor={#2}
       \pdfstringdef\__info_au: { #2 }
       \int_compare:nNnTF { \theau@ } < { 4 }
       { \xappto \infoauthors { \__info_au: , ~ } }
       {
         \int_compare:nNnTF { \theau@ } = { 4 }
         { \xappto \infoauthors { et~al. } }
         { }
       }
   }
\ExplSyntaxOff
\makeatother
% The byline badges already carry the iD links, so the class's separate
% "ORCID(s):" footnote would only repeat the same numbers.
\RenewDocumentCommand \printorcid { } { }
% Same author list printer as the class, except that the last two names are
% separated by a comma instead of "and" so the four-author line fits on one
% line together with the ORCID badges.
\ExplSyntaxOn
\RenewDocumentCommand \stmauthors { }
   {
     \group_begin:
     \stmAuthorSetup { type = authors }
     \l_stm_au_setup_tl
     \par \vskip\l_stm_augroup_before_dim
     \l_stm_augroup_align_tl
     \l_stm_augroup_size_tl
     \l_stm_augroup_shape_tl
     \l_stm_augroup_weight_tl
     \color{ \l_stm_augroup_color_tl }
     \bool_if:NTF \g_stm_augr_bool
       { \seq_use:Nn \g_stm_augr_seq { \par } }
       {
         \seq_use:cnnn { g_stm_au\int_use:N\g_stm_augr_int _seq }
         { ,~ } { ,~ } { ,~ }
       }
     \par\vskip\l_stm_augroup_after_dim
     \group_end:
   }
\ExplSyntaxOff

% Figures live in paper/figures/.
\graphicspath{{./figures/}{figures/}{paper/figures/}}

\begin{document}
\let\WriteBookmarks\relax

% Short title
\shorttitle{Tool-Result Errors in Language-Model Agents}

% Short author
\shortauthors{Rahaman et~al.}

% Main title of the paper
\title [mode = title]{Tool-Result Errors in Language-Model Agents: A Fixed-Prefix Measurement Study of Switches, Silent Propagation, and Recovery}

% ---------------------------------------------------------------------------
% Author block. The group, affiliations, ORCID iDs, and CRediT roles are reconciled to the
% byline above; any change before submission requires sign-off from all authors (order, spelling,
% e-mail addresses, ORCID iDs, CRediT allocation).
% ---------------------------------------------------------------------------

% First author and corresponding author
\author[1]{Md. Atikur Rahaman}[orcid=0009-0008-2945-4117]
\cormark[1]
\ead{mrahaman2310298@bscse.uiu.ac.bd}
\credit{Conceptualization, Data curation, Formal analysis, Investigation, Methodology, Project administration, Resources, Software, Validation, Visualization, Writing - original draft, Writing - review \& editing}



% Supervisor
\author[1]{Ohidujjaman Tuhin}[orcid=0000-0002-8776-7145]
\credit{Supervision, Validation, Writing - review \& editing}

% Address/affiliation
\affiliation[1]{organization={Computer Science and Engineering, United International University},
  city={Dhaka},
  citysep={},
  postcode={1212},
  country={Bangladesh}}

\cortext[cor1]{Corresponding author: Md. Atikur Rahaman}

% Abstract. Design numbers are sourced from README.md, ROADMAP.md,
% docs/PROTOCOL.md, and docs/DATA_CARD.md; confirmatory values from
% artifacts/analysis and artifacts/audit.
\begin{abstract}
  Tool-using language-model agents read tool outputs and act on them, but end-to-end benchmarks
  score the finished episode and cannot show whether an agent noticed a wrong price, date, or
  identifier. We report a frozen fixed-prefix measurement study of that moment: 800 real $\tau$-bench prefixes
  (400 airline, 400 retail), each ending after a tool observation, with the observation varied
  across four prose arms---faithful paraphrase, one non-critical change, one plausible critical
  error, and a critical error with a second contradiction. Three open-weight 8--9B models answer
  every arm greedily on an 18~GB H200 slice. The pre-registered content-sensitivity claim is not
  made: three gate criteria failed. The parser-failure rate is 1.55\% against the 1\% protocol
  threshold, the severity audit's $s2/s3$ decision-changing share is 88.3\% against 90\%, and crude
  recovery is 27.3/34.7/27.2\% by model against 50\% (Kaplan--Meier 32.4/47.6/28.2\% on the run
  risk set as a secondary with a competing-risk caveat). Within that bound, we report a measured shift among movers: the
  pooled conditional-logistic slope of the switch indicator on severity is 0.877 (95\% CI
  0.770--0.998), positive in all three models, with tier odds ratios against the prose placebo of
  1.27 (non-critical --- the $s1$ control is not inert, $p=.019$), 4.60 (critical), and
  11.21 (critical plus contradiction). These are assigned-severity effects: the audit finds
  7--16 of 60 critical-arm stimuli donor-unused (7 majority, 11 sheet-visible, 16 blind), and a
  bounded reclassification sensitivity steepens rather than flattens the $s1$--$s2$ contrast.
  Linearity is
  rejected, so the slope summarizes a threshold at $s=2$, not a per-level law. Of 1{,}964 critical-arm
  switches, at least 21.3\% silently repeat the injected value without flagging
  (recall-adjusted $\approx$28.5\%), and identifiers and money
  carry most of that propagation; status fields
  almost never do. Fewer than half of corrupted episodes return to the reference decision within
  four steps (22.0--44.6\% Kaplan--Meier secondary on the semantic risk set). Re-serializing the same observation without changing its
  values already shifts 20.3--45.6\% of G1 calibration decisions (main-run $s=0$ rates are
  19.8/31.9/38.1\%), so the slope is read over the prose arms
  $s0$--$s3$ only. Audits were blind-rated label-free (A-008) and checked by a multi-role agent panel,
  author-overseen. The G0 harvest audit is an AI pre-audit (40/40 adopted) with human verification
  deferred (D-021), so the freeze stands provisional on that item.
\end{abstract}

% Keywords. Each keyword is separated by \sep
\begin{keywords}
  Tool-result errors \sep LLM agents \sep Within-prefix design \sep Silent propagation \sep Recovery \sep $\tau$-bench
\end{keywords}

\begingroup
\hfuzz=120pt
\maketitle
\endgroup
\hypersetup{
  pdfauthor={Md. Atikur Rahaman, Ohidujjaman Tuhin},
  pdftitle={Tool-Result Errors in Language-Model Agents: A Fixed-Prefix Measurement Study of Switches, Silent Propagation, and Recovery}
}
\linenumbers

\section{Introduction}\label{sec:intro}

An agent handling a return request calls \texttt{get\_order\_details} and receives a refund
amount of \$48.50 instead of \$84.50. If the agent repeats \$48.50 to the customer, the episode
may still be recorded as completed. The failure is neither a wrong tool call nor a planning
error. It is a failure to notice that the observation conflicts with the policy or with the rest
of the context. A wrong flight date, a closed return window that the policy says is open, or a
cancelled reservation that still has seats produces the same shape of failure. Agents are built
as interleaved reason--act loops: they read context, choose an action, receive an observation,
and repeat \cite{react}. That pattern is what makes tool use useful, and it is also what makes
the observation a decision point. The next action is conditional on what the tool returned. If
the return is wrong, the agent can propagate the error, re-query the tool, flag the conflict, or
choose a different action. Aggregate task scores do not separate these outcomes.

End-to-end agent benchmarks score the final state of an episode. $\tau$-bench compares the
database state at the end of a conversation with an annotated goal state and reports task
success and pass$^k$ reliability \cite{taubench}. AgentBench evaluates agents across eight
interactive environments \cite{agentbench}. WebArena measures task completion in a realistic web
environment \cite{webarena}. OSWorld, WebVoyager, Mind2Web, GAIA, SWE-bench, and AppWorld extend
this style of end-to-end evaluation to desktop control, live web navigation, general assistance,
software repair, and interactive coding \cite{osworld,webvoyager,mind2web,gaia,swebench,appworld}.
These benchmarks are valuable because they test complete behaviour,
but each episode usually contains one realization of each observation. Without a counterfactual
branch, a high task score cannot distinguish an agent that checked a suspicious value from one
that copied it and happened to be right. Tool-use benchmarks focus on a different slice of the
problem. ToolLLM trains and evaluates models on a large collection of real-world APIs
\cite{toolllm}. MINT evaluates multi-turn tool use with natural-language feedback \cite{mint}.
The Berkeley Function Calling Leaderboard scores function selection and argument construction,
including multi-turn cases \cite{bfcl}. API-Bank, ToolSandbox, and the tool-learning survey cover
API reuse, stateful multi-turn tool use, and the broader design space
\cite{apibank,toolsandbox,tool_survey}. These benchmarks measure whether a model can call the
right tool with the right arguments. They do not hold a real observation fixed and vary only its
content.

Prior work shows that context and content influence model behaviour. Performance varies with the
position of relevant information in long contexts \cite{lostmiddle}. Models can agree with a
user's stated belief over a truthful answer, a behaviour known as sycophancy \cite{sycophancy}.
Tool-returned content can hijack an agent in prompt-injection attacks: AgentDojo evaluates
attacks and defences in tool-using environments \cite{agentdojo}, and InjecAgent benchmarks
indirect prompt injections in tool-integrated agents \cite{injecagent}. ToolEmu uses an emulated
sandbox to identify risky tool-use behaviour \cite{toolemu}. Instruction-hierarchy training and
spotlighting are two defenses against such content \cite{instruction_hierarchy,spotlighting}.
These studies establish that
observation content can change what an agent does. They study position, social pressure, or
adversarial instructions. The 2025--26 literature studies quiet tool failures directly.
ToolMaze varies tool perturbations along a $2\times2$ taxonomy and reports a perturbation
recovery rate, finding the sharpest drops under implicit semantic failures \cite{toolmaze}.
``Loud Failures, Quiet Failures'' injects four typed faults into multi-turn agents and
records whether the agent notices, changes plan, recovers, or repeats \cite{loudfailures}.
``Fabrication After Tool Failure'' shows agents asserting values their tools never returned,
with a one-sentence retrieval-status requirement cutting the rate \cite{fabrication}.
A robustness study of inconsistent and incomplete retail tool responses measures how
strongly a failure changes the next decision's action distribution \cite{robustagents}, and
``Measuring What Tool-Using LLM Agents Actually Do'' pairs a schema-valid wrong argument
across seven models and reports a causal error-propagation effect \cite{measuringwhat}.
ToolScan, ToolSword, and ToolBeHonest map tool-use error patterns, safety failures, and
hallucinations \cite{toolscan,toolsword,toolbehonest}; AgentAbstain pairs should-act with
should-abstain tasks \cite{agentabstain}; PALADIN trains recovery from tool failures
\cite{paladin}. Prompt formatting alone can move model behaviour sharply \cite{sclar}.
Table~\ref{tab:neighbors} positions this study against these neighbors.

\begin{table}[pos=htbp]
  \centering
  \caption{Nearest neighbors on quiet tool failures. Each row cites a verified
  arXiv record (publication status not verified); the last row is this study.
  No single prior study combines real benchmark prefixes, a graded
  content-only manipulation, a pre-registered within-prefix estimator, a
  silent/flag/action decomposition, and a survival-based recovery estimate.}
  \label{tab:neighbors}
  \footnotesize
  \setlength{\tabcolsep}{3pt}
  \begin{tabular}{>{\raggedright\arraybackslash}p{2.6cm}>{\raggedright\arraybackslash}p{4.4cm}>{\raggedright\arraybackslash}p{4.4cm}>{\raggedright\arraybackslash}p{3.4cm}}
    \toprule
    \textbf{Study} & \textbf{Perturbation} & \textbf{Paired contrast} & \textbf{Recovery / persistence} \\
    \midrule
    ToolMaze \cite{toolmaze} & $2\times2$ tool perturbations; implicit semantic & dynamic replanning paths & perturbation recovery rate \\
    Loud/Quiet Failures \cite{loudfailures} & four typed faults at a controlled point & notice / plan-change / repeat & recover-vs-repeat coding \\
    Fabrication \cite{fabrication} & unusable payload with \texttt{status:ok} & 1{,}024-item honesty probe & retrieval-status intervention \\
    Robust.\ Analysis \cite{robustagents} & inconsistent/incomplete retail returns & next-decision action shift & failure-type identification \\
    Measuring What \cite{measuringwhat} & one wrong tool argument, 7 models & paired error propagation & fabricated-use detection \\
    ToolScan \cite{toolscan} & seven error patterns & error taxonomy & --- \\
    ToolSword \cite{toolsword} & noisy misdirection, error conflicts & safety across 3 stages & --- \\
    ToolBeHonest \cite{toolbehonest} & tool hallucinations, depth + breadth & diagnostic benchmark & --- \\
    Sclar et al.\ \cite{sclar} & prompt-format variation & sensitivity to spurious features & --- \\
    AgentAbstain \cite{agentabstain} & should-act vs should-abstain & abstention rate & --- \\
    PALADIN \cite{paladin} & tool-failure cases & recovery training & self-correction rate \\
    \midrule
    This study & graded $s=0\ldots3$ content errors on frozen $\tau$-bench prefixes & pre-registered conditional-logit slope & 4-step KM/RMST survival \\
    \bottomrule
  \end{tabular}
\end{table}

The gap is therefore a combination, not a vacuum. To our knowledge, no prior study combines
a pre-registered graded content manipulation on real $\tau$-bench prefixes with a
silent-propagation decomposition and a survival-based recovery estimate. This
paper targets that combination with a paired design, a real-prefix dataset frozen before outcomes, and a
pre-registered analysis. The contributions are four. First, ToolResult-800 provides 800 real $\tau$-bench
prefixes with five observation arms, perturbation records, renderer version, and hashes; the
manifest is frozen and hashed. Second, the design measures a within-prefix shift in the next decision
across ordered severity arms---a measured shift among movers---rather than a between-episode
correlation. Third, it decomposes a decision change into action switch, detection, and silent
propagation, so that a fluent wrong answer can be separated from a caught one. Fourth, it
estimates a four-step recovery curve, which speaks to whether one verification step is enough. A
null or flat slope is a legitimate outcome and is reported as a measurement result at a known
precision.

A 60-prefix pilot tested whether the manipulation is runnable. It used synthetic WebArena-lite /
$\tau$-style prefixes, \texttt{qwen3.5:4b}, greedy decoding, and three collapsed action codes. The
pilot found a consult rate of 54/60, placebo jitter of 0/60, and an error-switch rate of
5/60 = 8.3\% (95\% CI [1.3, 15.3] pp; McNemar $b=5$, $c=0$, exact $p=.0625$). Those numbers show
that the manipulation runs and that error content was associated with a different next action in
5 of 60 prefixes. They do not estimate
a slope: the sample is small, the model is a single 4B model, the prefixes are synthetic, the
action coding is coarse, and there are no recovery data. The pilot is therefore a screening
signal, and Section~\ref{sec:results} keeps it separate from the confirmatory analysis. The
manuscript reports the completed study: the base matrix, recovery sequences, seed-repeat
shard, and blind audits are done, so the results below replace the earlier placeholders. Three
gate criteria failed and are reported as failures: the parser-failure rate (1.55\% against the 1\%
protocol threshold), the severity audit's $s2/s3$ decision-changing share (88.3\% against 90\%),
and recovery (crude 27.3/34.7/27.2\% against 50\%). Every quantity in the results section
traces to a ledger, a manifest, or an analysis artifact.

\section{Conventional Method}\label{sec:conventional}

\paragraph{Agent and tool-use benchmarks}.\

$\tau$-bench is the closest source of material because it ships historical trajectories with
tool calls and task databases \cite{taubench}. Its airline and retail domains contain the
deterministic tool outputs that this study replays. The benchmark reports task success and
pass$^k$ reliability over complete episodes. AgentBench spans operating systems, databases, and
knowledge graphs \cite{agentbench}; WebArena supplies a realistic web environment
\cite{webarena}; BFCL isolates function selection and argument construction \cite{bfcl}; ToolLLM
scales API coverage \cite{toolllm}; MINT adds multi-turn feedback \cite{mint}. Surveys document
the broader agent design space \cite{agent_survey}. Each measures a
capability that matters for agents. None of them holds one real observation fixed across
counterfactual variants of its content, so none separates noticing a wrong value from copying
it. Holistic evaluation and contamination audits further show that benchmark scores depend on
protocol and leakage \cite{helm,contamination}.

\paragraph{Observation content as a decision variable}.\

The observation is a message like any other, and message content is known to move model
behaviour. Lost-in-the-middle shows that the position of relevant text changes recall
\cite{lostmiddle}. Sycophancy shows that a stated belief can override a correct answer
\cite{sycophancy}. Hallucination and uncertainty work shows that models assert unsupported
content and that confidence is often miscalibrated
\cite{truthfulqa,halueval,selfcheckgpt,semantic_uncertainty,calibration,ovadia}. AgentDojo and
InjecAgent show that tool-returned text can redirect an agent
\cite{agentdojo,injecagent}, and ToolEmu shows that emulated tool environments can surface risky
behaviour before deployment \cite{toolemu}. These are existence results about adversarial or
social content. The present study asks a different question: given a realistic, non-adversarial
observation that the donor agent used, how much does a plausible factual error change the next
decision, and does the change scale with severity? The distinction matters for practice, because
a guardrail aimed at injections does not necessarily catch a wrong refund amount.

\paragraph{Counterfactual and paired designs}.\

The design is a within-unit contrast \cite{causal,rubin1974}. Because every arm shares the prefix,
task difficulty, domain, donor style, and conversation history are held constant, and the
comparison does not depend on between-task variation. The cost of this choice is that the
estimate is local to the frozen prefixes and to the renderer. The benefit is that a decision
switch can be attributed to the observation's content rather than to a change of task. The
verbatim arm is a separate anchor: it differs in format as well as content, so it is excluded
from the primary slope and used to measure format and decoding jitter.

\paragraph{Error propagation, detection, and recovery}.\
Silent propagation is the failure mode in which the false value appears in the next decision
without an explicit flag. It is invisible to a task-success metric when the episode still ends
in a plausible state. Detection and action switch are the two visible alternatives: the agent
flags the conflict, or it changes the tool or arguments. Recovery asks whether a later step
returns to the verbatim-observation decision or raises a flag. Work on self-correction and
reflection shows that iterative feedback can help when a reliable signal exists, but that
self-correction without external feedback is limited
\cite{selfrefine,reflexion,critic,selfcorrect_limit}. The recovery literature in agent
evaluation is thin, and most benchmarks censor the episode at its end. This study fixes a
four-step horizon, reports Kaplan--Meier and restricted mean survival time regardless of the
model fit, and treats the exponential half-life as estimable only when the gates allow it.

\paragraph{The conventional estimand}.\
An end-to-end agent benchmark scores a completed episode $\tau$ by its terminal environment state,
\begin{equation}
  S(\tau)=\mathbf{1}\{\,z_T(\tau)\in\mathcal{G}_\tau\,\},
  \qquad
  \theta_{\mathrm{e2e}}=\mathbb{E}_{\tau}\!\left[S(\tau)\right],
  \label{eq:e2e}
\end{equation}
where $z_T(\tau)$ is the state after the last turn and $\mathcal{G}_\tau$ is the annotated goal
set for that task. The all-trials reliability reported by $\tau$-bench is the same quantity
applied to $k$ independent trials,
\begin{equation}
  \mathrm{pass}^{k}=\mathbb{E}_{\tau}\!\left[\prod_{j=1}^{k}S(\tau_j)\right],
  \label{eq:passk}
\end{equation}
with $\tau_1,\dots,\tau_k$ independent runs of the same task \cite{taubench}. Both estimands are
defined on the episode. An agent that copies a wrong refund amount and still reaches a
goal-consistent state contributes $S=1$ exactly like an agent that notices the error; an agent
that notices it but fails for an unrelated reason contributes $S=0$. The aggregate therefore
cannot identify the effect of the observation's content, which is the quantity the next section
targets.

\section{Proposed Method}\label{sec:method}

\paragraph{Unit and donors}.\

The unit is a frozen message prefix that ends immediately after a real $\tau$-bench tool
observation. The prefix contains the system policy, the user turns, the donor assistant turns,
and all prior tool calls and results. The target model produces the next decision; it never
generates or repairs the prefix. Every arm is generated from the same prefix, so the contrast is
within prefix in the sense of a within-unit contrast \cite{causal}. The donors are the
public gpt-4o and claude-3.5-sonnet $\tau$-bench trajectories: gpt-4o airline, gpt-4o retail,
sonnet-35-new airline, and sonnet-35-new retail. Separating donors from targets avoids a circular
design in which the model that produced the history also supplies the outcome. The target models
are open-weight 8--9B families, so the prefixes are off-policy for them. The design fixes the
prefix across arms and stratifies by donor, but it does not make the history on-policy; this is a
limitation, not a fixable defect of the sample. Figure~\ref{fig:design} shows the design.

\begin{figure}[pos=htbp]
  \centering
  \figfile[0.95\linewidth]{figures/fig1.png}
  \caption{Fixed-prefix design: one frozen prefix, five observation arms, and greedy next-decision generation.}
  \label{fig:design}
\end{figure}

\paragraph{Notation and formal design}.\
Table~\ref{tab:notation} fixes the symbols used below. For prefix $p$, $x_p$ is the frozen message
sequence and $o_p$ is the real observation at its end. The perturbation operator produces the
four prose observations,
\begin{equation}
  o_p^{(s)}=\Pi_s(o_p),\qquad s\in\{0,1,2,3\},
  \label{eq:perturb}
\end{equation}
where $\Pi_0$ is a faithful paraphrase that preserves the value set, $\Pi_1$ changes one
non-critical leaf, $\Pi_2$ changes one decision-critical leaf to a plausible wrong value, and
$\Pi_3$ changes two leaves and introduces a contradiction. The renderer maps the prefix and the
arm observation to a prompt,
\begin{equation}
  R_a(o_p)=
  \begin{cases}
    \operatorname{JSON}(o_p), & a=V,\\[2pt]
    T_{\mathrm{prose}}\!\left(o_p^{(s)}\right), & a=s\in\{0,1,2,3\},
  \end{cases}
  \label{eq:render}
\end{equation}
and target model $m$ returns the next decision
\begin{equation}
  d_{p,m,a}=f_m\!\left(x_p,R_a(o_p)\right).
  \label{eq:decision}
\end{equation}
where the design holds $(x_p,o_p,m)$ fixed and varies only $a$, so the contrast is within prefix by
construction. Identification rests on two conditions: across the prose arms, the renderer changes
the observation's content but not the task, tools, history, or policy; and the $s=0$ arm preserves
the observation's value set, which is checked for every prefix (100\% equality). The verbatim arm
differs in format as well as content, so it supplies the reference signature and the format anchor
but is excluded from the slope. Greedy decoding makes a generation deterministic given its prompt
up to serving nondeterminism, which the seed-repeat shard measures rather than assumes. The
estimand is local to the frozen prefixes, the two donors, the pinned renderer, and the three target
models.

\begin{table}[pos=htbp]
  \centering
  \caption{Notation. Symbols are scoped to the equation where they are defined:
  $S$ is the episode-success indicator in Eq.~\ref{eq:e2e} and the hatted
  Kaplan--Meier survival estimate in Eq.~\ref{eq:km}; $\alpha$ is the nominal
  level (Eq.~\ref{eq:mde}) and Krippendorff's agreement (Eq.~\ref{eq:alpha});
  $\beta$ is the severity slope except for the quadratic coefficient $\xi$ and
  the fixed power in $z_{0.8}$; $p$ is a prefix except in $p$-values; $h$ is a
  recovery step (the Weibull hazard rate is $r_W$); $H$ is the horizon (the
  alternative hypothesis is spelled out). The score, variance, Weibull shape,
  subset, and intercept collisions are renamed ($T_g$, $W_g$, $\nu$, $I$,
  $a_{p,m}$).}
  \label{tab:notation}
  \small
  \begin{tabular}{>{\raggedright\arraybackslash}p{2.1cm}>{\raggedright\arraybackslash}p{11.6cm}}
    \toprule
    \textbf{Symbol} & \textbf{Meaning} \\
    \midrule
    $p$, $m$ & frozen prefix; target model \\
    $a$ & arm, $a\in\{V,0,1,2,3\}$; $s(a)$ is its severity, with $s(V)$ undefined \\
    $x_p$, $o_p$ & frozen message sequence; real tool observation at its end \\
    $o_p^{(s)}$, $\Pi_s$ & observation under perturbation $\Pi_s$ (Eq.~\ref{eq:perturb}) \\
    $R_a$, $d_{p,m,a}$ & arm renderer ($T_{\mathrm{prose}}$ for prose arms); next decision returned by model $m$ (Eqs.~\ref{eq:render}--\ref{eq:decision}, $f_m$ the target model) \\
    $\sigma(\cdot)$ & canonical decision signature (tool and arguments, or action verb and slots) \\
    $Y$, $Y^{A}$, $Y^{F}$ & switch; action switch; explicit flag \\
    $Y^{U}$, $Y^{O}$ & silent propagation; other switch \\
    $C^{A},C^{F},C^{U},C^{O}$ & exclusive typology classes partitioning $Y$ (Eq.~\ref{eq:partition}) \\
    $a_{p,m}$, $\beta$ & prefix-by-model intercept; severity slope, log-odds per level (Eq.~\ref{eq:primary}) \\
    $\gamma_k$, $\xi$ & categorical-severity coefficients; quadratic coefficient (Eq.~\ref{eq:categorical}) \\
    $\Lambda$ & likelihood-ratio statistics $\Lambda_{\mathrm{cat}}$, $\Lambda_{\mathrm{quad}}$ (Eq.~\ref{eq:lrt}) \\
    $T_g$, $W_g$, $E_g$ & conditional score, variance, and mean for group $g$ (Eqs.~\ref{eq:score}--\ref{eq:sandwich}) \\
    $h$, $H$ & recovery step; recovery horizon ($H=4$) \\
    $\mu$, $t_{1/2}$ & per-step hazard; half-life $\ln2/\mu$ (Eq.~\ref{eq:recovery}) \\
    $\nu$, $\lambda$ & Weibull shape and scale; $S_W$, $r_W$ its survival and hazard rate (Eq.~\ref{eq:weibull}) \\
    $\hat S_{\mathrm{KM}}(h)$, RMST & Kaplan--Meier survival estimate; restricted mean survival (Eq.~\ref{eq:km}) \\
    $\tau$, $\theta_{\mathrm{e2e}}$ & an episode; the conventional end-to-end estimand (Eq.~\ref{eq:e2e}) \\
    $v^{\ast}$, $\mathcal{V}(d)$, $F$ & injected false value; typed values a decision asserts; flag indicator (Eqs.~\ref{eq:switch}--\ref{eq:silent}) \\
    $n_a$ & labeled rows in arm $a$ \\
    $p_o$, $p_e$ & observed and chance agreement (Eq.~\ref{eq:kappa}) \\
    \bottomrule
  \end{tabular}
\end{table}

\subsection{Sampling frame and eligibility}

\label{sec:supply}

The supply is measured, not assumed. The pinned frame is $\tau$-bench commit \texttt{59a200c},
with four trajectory files, six task databases, and a harness digest recorded in
\texttt{data/manifests/source\_pin.json}. Parsing every tool message followed by an assistant
message yielded 1{,}980 trajectories and 13{,}605 post-tool candidates, of which 11{,}318 carry
JSON-parseable observations. Excluding the profile and PII tools
(\texttt{get\_user\_details}, \texttt{find\_user\_id\_by\_name\_zip},
\texttt{find\_user\_id\_by\_email}, and \texttt{modify\_user\_address}) removed 3{,}692
candidates. The operational value-used rule, which keeps a prefix only when at least one
observation scalar appears in the donor's factual continuation after normalization, left
5{,}516 eligible prefixes: 1{,}434 airline and 4{,}082 retail. Replay against the pinned task
database matched 5{,}511 of 5{,}516 eligible prefixes (99.91\%). Table~\ref{tab:supply} reports
the funnel. A prefix is eligible if it contains a real tool message, a reconstructable prefix, at least one
scalar value used in the donor's factual continuation, PII-scrubbed content, a replay match of at
least 98\%, and a lossless $s=0$ rendering. Replay is the step that makes an observation a real
environment output: the harness loads the pinned task database, executes the prefix's tool calls
from the initial state, and compares the final observation with the historical one using a hash
and normalized values. The 98\% threshold is not relaxed. Full-pool replay was 87.39\%, and the
gap is the historical \texttt{get\_user\_details} \texttt{province}/\texttt{state} schema drift,
which is disclosed in \texttt{docs/DECISIONS.md}; the replay population is the eligible
non-profile pool.

\begin{table}[pos=htbp]
  \centering
  \caption{Sampling frame and selected sample from the pinned $\tau$-bench frame.}
  \label{tab:supply}
  \small
  \begin{tabular}{>{\raggedright\arraybackslash}p{6.6cm}rr}
    \toprule
    \textbf{Stage} & \textbf{Count} & \textbf{Note} \\
    \midrule
    Source trajectories                 & 1{,}980 & four donor files \\
    Post-tool candidates                & 13{,}605 & tool message + assistant reply \\
    JSON-parseable observations         & 11{,}318 & \\
    Profile/PII-tool candidates excluded & 3{,}692 & four excluded tools \\
    Eligible value-used                 & 5{,}516 & airline 1{,}434; retail 4{,}082 \\
    Replay-matched (eligible pool)      & 5{,}511 & 99.91\% \\
    \quad full-pool replay              & 87.39\% & schema drift, D-011 \\
    Cap-10 pool                         & 1{,}586 & airline 442; retail 1{,}144 \\
    Selected (primary + reserve)        & 880 & seed 20261108 \\
    \quad primary                       & 800 & 400 airline; 400 retail \\
    \quad reserve                       & 80 & 40 airline; 40 retail \\
    Calibration prefixes                & 60 & 30 airline; 30 retail; disjoint \\
    Rendered arms                       & 4{,}400 & 880 $\times$ 5 \\
    Perturbation records                & 2{,}640 & 880 each for $s=1,2,3$ \\
    \bottomrule
  \end{tabular}
\end{table}

Sampling is stratified by domain, donor, step-index bucket, and tool family, with seed 20261108.
The study selected 880 prefixes: 800 primary (400 airline, 400 retail) plus 80 reserve (40
airline, 40 retail). The sample is frozen before any target-model output is generated. The
manifest records prefix IDs, message hashes, observation hashes, arm assignments,
perturbation-record hashes, the renderer version, the field-registry hash, and the preregistration
hash. Every analysis table must trace to a manifest row and a ledger row. The airline arm is the
binding supply constraint: the cap-10 pool has 442 airline prefixes, so 400 primary plus 40
reserve leaves a margin of two before attrition.

\paragraph{Arms and rendering}.\ \label{sec:arms}
All arms use a real observation from the same prefix. V is the exact observation as raw JSON. The
four prose arms use one deterministic key--value/bullet template; only content changes. No
LLM-generated text enters the manipulation. A script flattened the V value set and the $s=0$ value
set and verified equality; any loss excluded the prefix. Table~\ref{tab:arms} defines the arms
and Figure~\ref{fig:arms} shows the construction for one prefix.

\begin{table}[pos=htbp]
  \centering
  \caption{Observation arms and their roles in the slope.}
  \label{tab:arms}
  \small
  \begin{tabular}{>{\raggedright\arraybackslash}p{1.3cm}c>{\raggedright\arraybackslash}p{6.3cm}>{\raggedright\arraybackslash}p{4.6cm}c}
    \toprule
    \textbf{Arm} & \textbf{Level} & \textbf{What changes} & \textbf{Role} & \textbf{In slope} \\
    \midrule
    V     & anchor & Exact real observation, raw JSON. & Format/jitter anchor & no \\
    $s=0$ & 0      & Faithful prose; all values preserved. & Placebo & yes \\
    $s=1$ & 1      & One non-critical field changed or omitted. & Negative control & yes \\
    $s=2$ & 2      & One critical field changed to a plausible wrong value. & Dose & yes \\
    $s=3$ & 3      & Critical wrong value plus a second detectable contradiction. & Dose (compound) & yes \\
    \bottomrule
  \end{tabular}
\end{table}

\begin{figure}[pos=htbp]
  \centering
  \figfile[0.95\linewidth]{figures/fig2.png}
  \caption{Arm construction illustrated for one prefix.}
  \label{fig:arms}
\end{figure}

\subsection{Perturbation records and the severity rubric}



The construction rules are pre-specified. For $s=1$, the perturbed field must be absent from the
donor's factual continuation; if a non-critical change is associated with as many switches as $s=2$, the
field was misclassified or the model is hypersensitive, and both are audit findings. For $s=2$,
the wrong value must be plausible for the entity and must change the correct action if trusted;
absurd values are rejected. For $s=3$, the critical error is combined with a second inconsistency
that a careful reader can detect, such as ``cancelled'' while seats are available or
``non-refundable'' inside the refund window. Calibration checks that $s=3$ is not a ceiling.
Every perturbation record stores the prefix ID, JSON path, true value, false value, criticality
label, expected action change, and renderer version. PII fields are never perturbed.

A worked example makes the coding concrete. A user asks to return headphones from order W1842.
The tool returns status ``delivered'', delivery three days ago, a 30-day return window that is
open, a refund of \$84.50, and a product rating of 4.4. Under $s=2$, the refund becomes \$48.50.
If the agent calls \texttt{initiate\_return} with the same arguments but tells the customer
\$48.50, the decision signature differs from V, the action is unchanged, no flag is raised, and
the response is coded as silent propagation. Under $s=3$, the refund is \$48.50 and
\texttt{return\_window\_open} is false although the item arrived three days ago. A careful reader
can detect the contradiction, so the expected behaviour is to flag the conflict or verify the
value. Severity is ordinal, not interval. The linear slope is a summary of the graded contrast;
the $s=2$ versus $s=3$ contrast is always reported separately because $s=3$ is intentionally
compound. A quadratic term and a categorical-severity likelihood-ratio test are pre-registered
diagnostics.

\paragraph{Audits, blinding, and PII}.\

Raw $\tau$-bench trajectories may contain names, emails, addresses, phone numbers, and payment
identifiers. Profile tools are excluded, direct identifiers become stable synthetic tokens, and a
prefix is excluded if decision-relevant content is itself PII. A 40-item audit checks
eligibility, PII, relevance, and replay before freeze; in this run the sheet was filled by an
AI-assisted mechanical pre-audit (D-021), and the human verification required by the
preregistration is reported as a deviation rather than as passed. A separate re-scan runs before
release.
Raw trajectories, weights, caches, credentials, and run logs never enter the repository.
Automatic labels are primary and are validated by blind audits. The severity audit uses 120
observations (30 per level) rated by three separate blind rater sessions; the outcome audit
uses 160 responses rated by two; the recovery audit uses 80 sequences rated by two from per-step
evidence, where each step's parsed decision and flag bit are given. The reported round is
label-free: the packets exclude the mechanical reference labels and every automatic answer, and
each rater sees only the evidence. The severity packet shows both changed leaves for compound
perturbations, in shuffled item order, and excludes arm-derived annotations; the outcome packet
includes the injected true/false values; the recovery packet includes the per-step evidence.

A first round was label-visible; the author accepted it under
D-035, and it is archived as a sensitivity. The label-free round is primary (A-008). Its sheets were
checked by a multi-role agent panel, author-overseen (one verifier per audit, plus a second-rater
pair on a stratified 120-item subset); the panel is recorded under D-043. Its ratings are agent-produced and share one model
family, so they establish reproducibility rather than human independence.
The ratings were produced by blind rater sessions operated with an AI agent. The severity gate
requires Krippendorff's $\alpha \ge .75$ \cite{krippendorff}, a Jonckheere--Terpstra trend test at
$p<.01$ \cite{jonckheere1954}, at least 90\% of $s=1$ items rated non-decision-changing, at least
90\% of $s=2$ and $s=3$ items rated decision-changing, and no more than 20\% of $s=3$ items rated
absurd. The outcome gate requires Cohen's $\kappa \ge .80$ \cite{cohen1960}, silent-propagation
positive predictive value at least .90, parser failure below 1\%, and seed-repeat stability at
least .90. The recovery gate requires $\kappa \ge .75$. If outcome $\kappa$ falls below .70 after
one rubric revision, the study downgrades to action-level $Y^{A}$ only and does not make the
content-sensitivity claim.



Figure~\ref{fig:pipeline} shows how the stages connect, from the pinned upstream checkout to the
ledger row that carries each reported number.

\subsection{Next-decision prompt}



The prompt is version 3 (decision D-020) and is frozen before calibration. It renders the call
history as plain text and ends with one instruction turn that asks for the next decision only. The model must
return JSON: either a tool call with a tool name and arguments, or a final answer with an action
verb and decision-critical slots. The prompt does not mention that the observation may be
perturbed, and it does not ask the model to verify values. That omission is deliberate, because
the outcome of interest is what the model does by default. A prompt that warned about errors
would measure compliance with an instruction rather than sensitivity to content. The prompt
version is frozen at G0 (prompt v3) and its hash is recorded in the manifest. The parser was
revised during calibration (D-024, D-025) and once after the blind audit (D-028); the final
parser is \texttt{toolresult-parser-v4}, recorded in every ledger row. The manifest records
config hashes and the renderer version, not a parser field.

\paragraph{Parser and canonicalization}.\

The parser reads the JSON decision and canonicalizes it before comparison. Canonicalization sorts
keys, normalizes numbers, converts dates to ISO form, lower-cases entities, and collapses
whitespace and punctuation. A decision signature is the tool name plus the canonicalized
arguments, or the final answer's normalized action verb plus decision-critical slots. A
generation with a parser failure is retained in the ledger and reported; it is not silently
dropped. A parser-error rate above 5\% in a shard fails fast rather than burning GPU hours, and
the overall parser failure must stay below 1\% for G3.

\paragraph{Context and length envelope}.\
The longest observed prompt is 20{,}681 tokens, and the decision budget is 1{,}024 tokens. The
study therefore adopts a 24{,}576-token context. The choice was audited in D-019: two synthetic
20k-token prompts pass on all three models. No prefix is truncated and no observation is
trimmed; a prefix whose prompt did not fit would be excluded before freeze, but none does. The
24k context is also the constraint that drives the batch-size fit check in
Section~\ref{sec:experiment}.

\paragraph{Serving stack and the model trio}.\

All three models are 4-bit AWQ quantized \cite{awq} and served one at a time with vLLM 0.33.0
\cite{vllm} and torch 2.14.0 on the approved H200 MIG 1g.18gb slice: 18~GB of VRAM, driver
580.178.04, CUDA 13.5.
Decoding is greedy with a fixed seed. Table~\ref{tab:models} lists the trio. Qwen3.5-9B runs with
its thinking block disabled, because a reasoning block would change the decision format and the
token budget. If Llama-3.1 access fails, the pre-registered substitutes are InternLM2.5-7B-Chat
or Yi-1.5-9B-Chat, recorded before calibration and never selected after seeing target-model
outcomes. The local RTX 3060 (6~GB) is used only for imports, smoke tests, renderer checks, and
tiny forward passes; it cannot host an 8--9B AWQ model at a 24k context.

\begin{table}[pos=htbp]
  \centering
  \caption{Target model trio served under greedy decoding (seed 20261108). The served
  per-model concurrent-sequence caps are the calibration outcome (D-016); exact weight
  revisions are pinned in \texttt{configs/models.yaml} and the Llama-3.1 base weights are
  described in \cite{llamaherd}. Non-study fit checks on synthetic prompts pass at 24k
  context for all three (batch 4), so no prefix is truncated.}
  \label{tab:models}
  \small
  \begin{tabular}{>{\raggedright\arraybackslash}p{4.9cm}cc>{\raggedright\arraybackslash}p{2.0cm}>{\raggedright\arraybackslash}p{3.9cm}}
    \toprule
    \textbf{Model} & \textbf{Params} & \textbf{Served seqs} & \textbf{Quantization} & \textbf{Notes} \\
    \midrule
    Qwen3.5-9B-AWQ                        & 9B & 8 & AWQ 4-bit & thinking block disabled \\
    Ministral-3-8B-Instruct-2512-AWQ      & 8B & 4 & AWQ 4-bit & instruction-tuned \\
    Llama-3.1-8B-Instruct-AWQ-INT4        & 8B & 4 & AWQ 4-bit & access fallback: InternLM2.5-7B-Chat or Yi-1.5-9B-Chat \\
    \bottomrule
  \end{tabular}
\end{table}

\begin{figure}[pos=htbp]
  \centering
  \figfile[0.95\linewidth]{figures/fig3.png}
  \caption{End-to-end pipeline from pinned checkout to ledger row.}
  \label{fig:pipeline}
\end{figure}



\subsection{Decision signature and labels}



For prefix $p$, model $m$, and severity $s \in \{0,1,2,3\}$, let $d=d_{p,m,s}$ and
$d_V=d_{p,m,V}$ be the decisions under arm $s$ and the verbatim arm, and let
$F_{p,m,s}=\mathbf{1}[\text{flag detector fires}]$. The switch indicator is
\begin{equation}
  Y_{p,m,s}
  = \mathbf{1}\!\left[\sigma(d)\neq\sigma(d_V)\right]\ \lor\ F_{p,m,s},
  \label{eq:switch}
\end{equation}
where $\sigma(\cdot)$ is the canonical decision signature. A flag alone therefore counts as a
switch even when the signature is unchanged, because an explicit doubt response is a
decision-relevant deviation. The component indicators are
\begin{align}
  Y^{A} &= \mathbf{1}\!\left[\operatorname{act}(d)\neq\operatorname{act}(d_V)\right],
  \label{eq:action}\\
  Y^{F} &= F_{p,m,s},
  \label{eq:flag}\\
  Y^{U} &= Y\,(1-Y^{A})(1-Y^{F})\,
           \mathbf{1}\!\left[v^{\ast}\in\mathcal{V}(d)\right]\,
           \mathbf{1}\!\left[v^{\ast}\notin\mathcal{V}(d_V)\right],
  \label{eq:silent}\\
  Y^{O} &= Y\,(1-Y^{A})(1-Y^{F})(1-Y^{U}),
  \label{eq:other}
\end{align}
where $\operatorname{act}(d)$ is the tool name for a tool call and the action verb for a final
answer, $\mathcal{V}(d)$ is the set of typed values the decision asserts (tool arguments, or the
status, money, date, time, and identifier slots of a final answer), and $v^{\ast}$ is the injected
false value. $Y^{U}$ is defined only for $s\ge2$; $Y^{O}$ collects signature changes not
attributable to an action switch, a flag, or the injected value.
The raw indicators are not mutually exclusive, because an action switch can also carry a flag.
The reported typology uses the exclusive indicators
\begin{equation}
  C^{A}=Y^{A},\quad
  C^{F}=(1-Y^{A})Y^{F},\quad
  C^{U}=(1-Y^{A})(1-Y^{F})Y^{U},\quad
  C^{O}=(1-Y^{A})(1-Y^{F})(1-Y^{U})Y,
  \label{eq:partition}
\end{equation}
which satisfy $Y=C^{A}+C^{F}+C^{U}+C^{O}$ and give the action-switch, flag-only,
silent-propagation, and other shares reported in Section~\ref{sec:results}. The strict form makes
the silent-propagation requirement explicit: a model that re-queries the tool and also repeats the
wrong value is not the same failure as one that quietly repeats it. Arm-level rates are the
corresponding sample means, for example
$\hat\pi_a=n_a^{-1}\sum_{p,m}Y_{p,m,a}$ with $n_a$ the number of labeled rows in arm $a$. The
strict and lenient sensitivity analyses replace $Y$ with, respectively, a signature difference
that ignores flags and an action-only switch; the slope is refitted under each coding. The
strict coding is the signature-only $Y$ of the preregistered outcome definition (flags
excluded); because it leaves the slope unchanged (0.876 vs 0.877) and the flag/other
components ($Y^{F}$, $Y^{O}$) are decomposed separately, the outcome-definition difference is
handled by this reported sensitivity, with no prereg amendment
(\texttt{artifacts/analysis/inference\_robustness.json};
\texttt{prereg/DEPOSIT\_NOTE.md}).

\paragraph{Primary estimand}.\

The primary model is a conditional (fixed-effects) logistic regression stratified by prefix:
\begin{equation}
  \operatorname{logit}\Pr(Y_{p,m,s}=1) = a_{p,m} + \beta_m s,
  \qquad s \in \{0,1,2,3\},
  \label{eq:primary}
\end{equation}
where $a_{p,m}$ is a prefix-by-model intercept and $\beta_m$ is the log-odds change per
one-level increase in severity. The pooled model stratifies by prefix $\times$ model and clusters
standard errors by prefix. V is excluded from the slope because it differs in format as well as
content; $s=0$ is the placebo that absorbs format jitter. The primary test is one-sided
for the alternative $\beta>0$ at $\alpha=.05$, reported two-sided for completeness. H1 is supported only if the
pooled slope is positive and the per-model slope is positive in at least two of three models with
the same sign. The second condition is a replication gate, not a model ranking.

\paragraph{Conditional likelihood and inference}.\
Because $a_{p,m}$ is a nuisance intercept, $\beta$ is estimated by conditional maximum
likelihood \cite{chamberlain1980}. For group $g=(p,m)$, let $k_g=\sum_{i\in g}y_i$ be the number of switches and let
$s_i$ be the severity of row $i$. The group contribution is
\begin{equation}
  L_g(\beta)=
  \frac{\exp\!\left(\beta\sum_{i\in g}s_i y_i\right)}
       {\sum_{I\subseteq g,\,|I|=k_g}\exp\!\left(\beta\sum_{i\in I}s_i\right)},
  \label{eq:clogit}
\end{equation}
and $\hat\beta$ maximizes $\ell(\beta)=\sum_g\log L_g(\beta)$. Groups with $k_g=0$ or
$k_g=n_g$ are uninformative and drop out; with at most four rows per group, the denominator
enumerates at most $\binom{4}{2}=6$ subsets. The score and conditional variance are
\begin{equation}
  T_g(\beta)=\sum_{i\in g}s_i y_i-E_g(\beta),
  \qquad
  W_g(\beta)=\operatorname{Var}_g\!\left(\sum_{i\in I}s_i\right),
  \label{eq:score}
\end{equation}
where $E_g$ and the variance are taken under the conditional distribution in
Equation~\ref{eq:clogit}. The prefix-cluster sandwich variance is
\begin{equation}
  \widehat{\operatorname{Var}}(\hat\beta)=
  \frac{\sum_{p}\left(\sum_{g:\,p(g)=p}T_g(\hat\beta)\right)^2}
       {\left(\sum_{g}W_g(\hat\beta)\right)^2},
  \label{eq:sandwich}
\end{equation}
which is the clustered standard error reported in the tables \cite{cameron2015}; the model-based
variance $1/\sum_g W_g$ is retained only as a diagnostic. Confidence intervals use a prefix-cluster
percentile bootstrap: resample prefixes with replacement, refit $\beta$ in each replicate, and
take the 2.5th and 97.5th percentiles over 2{,}000 replicates (seed 20261108) \cite{efron1979}.

\paragraph{Linearity and multiplicity}.\

The slope assumes equal spacing of $s=0,\dots,3$. The categorical-severity model is
\begin{equation}
  \operatorname{logit}\Pr(Y_{p,m,s}=1)
  =a_{p,m}+\sum_{k=1}^{3}\gamma_k\mathbf{1}[s=k],
  \label{eq:categorical}
\end{equation}
with $s=0$ as the reference, and the quadratic model adds a term $\xi s^2$ to
Equation~\ref{eq:primary}. The two pre-registered linearity diagnostics are likelihood-ratio
tests,
\begin{equation}
  \Lambda_{\mathrm{cat}}=2(\ell_{\mathrm{cat}}-\ell_{\mathrm{lin}})\sim\chi^{2}_{2},
  \qquad
  \Lambda_{\mathrm{quad}}=2(\ell_{\mathrm{quad}}-\ell_{\mathrm{lin}})\sim\chi^{2}_{1},
  \label{eq:lrt}
\end{equation}
and the tier-specific odds ratios are $\widehat{\mathrm{OR}}_k=e^{\hat\gamma_k}$ for
$k=1,2,3$. If linearity is rejected, the paper reports these tier estimates and treats the slope
as a summary rather than a law. The fixed testing sequence is: pooled $\beta$; per-model $\beta$
with Holm correction across the three models \cite{holm1979}; descriptive silent propagation;
recovery; exploratory heterogeneity. No post-hoc subgroup claim is promoted to confirmatory
status, and no model ranking is a confirmatory claim. The paired $s=2$ versus $s=3$ contrast uses
the exact McNemar test. With $b$ pairs that switch at $s=3$ but not $s=2$ and $c$ pairs that
switch at $s=2$ but not $s=3$, the two-sided $p$-value is
\begin{equation}
  p_{\mathrm{McNemar}}=
  2\sum_{j=0}^{\min(b,c)}\binom{b+c}{j}2^{-(b+c)},
  \label{eq:mcnemar}
\end{equation}
and the one-sided version sums only to $b$ \cite{mcnemar1947}. The per-model
exact McNemar tests are the primary paired contrast: within one model each prefix contributes
exactly one $s=2$-vs-$s=3$ pair, so prefixes are the independent units and no cross-model
clustering applies. The pooled exact test over 2{,}251 model-prefix pairs is secondary, because
up to three pairs share one prefix and the exact $p=2.1\times10^{-20}$ is anti-conservative for
prefix-level inference. Two cluster-robust checks bound the concern: a model-stratified
(CMH-style) McNemar agrees ($\chi^2=81.67$, $p=1.6\times10^{-19}$), and a design-effect grid
with at most $m=3$ pairs per prefix keeps the pooled contrast significant at every correlation
($\rho=0.5$: $\chi^2=40.83$, $p=1.7\times10^{-10}$;
\texttt{artifacts/analysis/inference\_robustness.json}). The severity audit's monotonicity
test is Jonckheere--Terpstra \cite{jonckheere1954},
\begin{equation}
  J=\sum_{k<l}\ \sum_{x\in G_k}\ \sum_{y\in G_l}
  \left(\mathbf{1}[y>x]+\tfrac12\mathbf{1}[y=x]\right),
  \label{eq:jt}
\end{equation}
over the ordered severity groups $G_0,\dots,G_3$, with a Monte-Carlo permutation $p$-value from
20{,}000 relabelings (seed 20261108); half credit for ties keeps the statistic conservative.
Donor-specific and task-specific slopes and the exploratory model $\times$ severity and task
$\times$ severity interactions are reported separately. Bootstrap intervals use 2{,}000
replicates clustered by prefix.

\paragraph{Recovery}.\

Recovery uses the risk set of prefixes whose decision differed from the reference at the injected step ($Y^{U}=1$ or
$Y^{A}=1$) and did not flag there, separately for $s=2$ and $s=3$. Prefixes that flagged at
injection are coded $h=0$ for detection and are not at risk for later recovery. The episode
continues for $H=4$ steps against the real environment. Recovery at step $h$ is the first step
where the signature matches V (R1: reference decision) or the model explicitly flags (R2:
detection); a flag at the same step is coded R2 before the signature check. A sequence that does
neither by $H=4$ is censored. The model is discrete-time
exponential with interval censoring (Equation~\ref{eq:recovery}),
\begin{equation}
  \Pr(\text{recover at } h \mid \text{not yet}) = 1-e^{-\mu},
  \qquad
  t_{1/2} = \frac{\ln 2}{\mu},
  \label{eq:recovery}
\end{equation}
over $H=4$. The Kaplan--Meier estimator is
\begin{equation}
  \hat S_{\mathrm{KM}}(h) = \prod_{i=1}^{h}\left(1-\frac{d_i}{n_i}\right),\qquad
  \mathrm{RMST} = \sum_{h=1}^{H} \hat S_{\mathrm{KM}}(h),
  \label{eq:km}
\end{equation}
with $d_i$ events and $n_i$ at risk at step $i$. Kaplan--Meier and RMST are reported regardless
of the fit \cite{kaplan1958}. The parametric check is a Weibull model with survival and hazard
\begin{equation}
  S_{W}(t)=\exp\!\left[-\left(\frac{t}{\lambda}\right)^{\nu}\right],
  \qquad
  r_{W}(t)=\frac{\nu}{\lambda}\left(\frac{t}{\lambda}\right)^{\nu-1},
  \label{eq:weibull}
\end{equation}
where $\nu$ is the shape and $\lambda$ the scale. The exponential is the special case
$\nu=1$; a likelihood-ratio test of $\nu=1$ at $p<.01$ checks the assumption. If the shape
rejects, or if fewer than 50\% of the risk set recovers by $H=4$,
$t_{1/2}$ is labelled exploratory and the curve and RMST are the result. Transient injection is
the primary regime. The persistent regime was designed to re-apply the false value to each new
observation on 25\% of prefixes; in execution the corrupted path recurred rarely, so the step
prompts are byte-identical to the transient regime for 97--100\% of shared steps and the regime
contrast is reported as uninformative (Section~\ref{sec:results}). Algorithm~\ref{alg:recovery}
gives the endpoint coding.

\begin{algorithm}[htbp]
  \caption{Recovery endpoint coding for one at-risk sequence.}
  \label{alg:recovery}
  \begin{algorithmic}[1]
    \Require prefix $p$ in the risk set, reference decision $d_V$, horizon $H=4$
    \For{$h=1,\dots,H$}
      \State generate $d_h$ and parse it
      \If{$d_h$ is not parsed}
        \State \Return censored at $h-1$ \Comment{parse break ends the sequence}
      \EndIf
      \If{the flag detector fires on $d_h$}
        \State \Return R2 at $h$ \Comment{detection; checked before the signature}
      \EndIf
      \If{$\sigma(d_h)=\sigma(d_V)$}
        \State \Return R1 at $h$ \Comment{reference decision}
      \EndIf
      \If{$d_h$ is a final answer}
        \State \Return censored at $h$
      \EndIf
      \State append $d_h$; invoke its tool; if persistent, re-apply the injected value;
      append the observation
    \EndFor
    \State \Return censored at $H$
  \end{algorithmic}
\end{algorithm}

\paragraph{Power}.\

The design assumes a 4\% placebo switch rate and a random prefix intercept with standard
deviation 1.0. The minimum detectable effect follows Equation~\ref{eq:mde}. At $N=800$ the per-model minimum detectable effect is
0.203 logit per level (odds
ratio 1.22) at 80\% power and $\alpha=.05$ two-sided; the pooled minimum detectable effect is at
most 0.15. The observed $s0$ switch rate is 30.1\%, far above the 4\% design assumption, so the
design MDE is not the realized precision: at the analyzed sample and observed cluster SE the
realized pooled MDE is 0.166 logit/level (OR 1.18, 80\% power), and realized power at the observed
effect is 1.0. The pilot-implied value of about 0.8 is roughly four times the per-model MDE, so power
is not the binding constraint. Construct validity is. The recovery half-life has wider
uncertainty: with 160--240 $s=3$ switchers, the per-model interval is about $\pm16$--20\% and the
pooled interval about $\pm9$--11\%, while between-model and between-severity differences need
32--50\%. That is why recovery is an estimation target and not a ranking.
Figure~\ref{fig:analysis} summarises the analysis path from the switch indicator to the recovery
model.

\begin{figure}[pos=htbp]
  \centering
  \figfile[0.95\linewidth]{figures/fig4.png}
  \caption{Conceptual schematic of the analysis path (not study data): the within-prefix
  switch indicator, the conditional-logit slope, the silent-propagation decomposition, and
  the four-step recovery model with Kaplan--Meier and RMST reporting (Figure~\ref{fig:design}
  shows the design).}
  \label{fig:analysis}
\end{figure}

\section{Experiment}\label{sec:experiment}

\paragraph{Why sequential gates}.\

Gates are sequential and frozen before the main run. A gate cannot be skipped, relaxed after
seeing data, or re-run on a different sample without a logged amendment. Each gate report states
the criterion, the observed value, the pass/fail decision, and the consequence. A failed gate
changes the claim, not the data. Table~\ref{tab:gates} lists the criteria and the fail action for
G0 through G4. G0 guards supply and freeze integrity, G1 the calibration harness, G2 the compute
budget, and G3 and G4 the slope and recovery claims.
\begin{table}[pos=htbp]
  \centering
  \caption{Sequential gates and stop rules.}
  \label{tab:gates}
  \footnotesize
  \setlength{\tabcolsep}{4pt}
  \begin{tabular}{>{\raggedright\arraybackslash}p{0.9cm}>{\raggedright\arraybackslash}p{9.6cm}>{\raggedright\arraybackslash}p{4.2cm}}
    \toprule
    \textbf{Gate} & \textbf{Pass criterion} & \textbf{If it fails} \\
    \midrule
    G0 & At least 880 eligible; at least 440 per domain before audit and 400 after audit; replay at least 98\%; 40-item audit at least 36/40; freeze hash stable. & Fall back to 500 prefixes and 2 models, or stop. \\
    G1 & Per model: parse at least 95\%; placebo jitter at most 20\%; $s=3$ switch 10--80\%; format refusal below 5\%. & Fix prompt or parser and rerun; no main-run start. \\
    G2 & Stop at 45 GPU-hours or 200~GB storage, whichever comes first; ledger complete. & Analyse completed shards; disclose the shortfall. \\
    G3 & Pooled $\beta>0$; positive in at least 2 of 3 models with the same sign; severity $\alpha \ge .75$; outcome $\kappa \ge .80$ (or adjudicated at least .70 with downgrade); recovery $\kappa \ge .75$; silent-propagation PPV at least .90; parser failure below 1\% (protocol criterion from \texttt{docs/GATES.md}, absent from the frozen prereg G3 row; A-011). & Measurement paper; no causal or content-sensitivity claim. \\
    G4 & At least 50\% of the $s=3$ risk set recovers by $H=4$; Weibull shape not rejected at $p \ge .01$. & Report RMST and the curve; $t_{1/2}$ exploratory only. \\
    \bottomrule
  \end{tabular}
\end{table}
\begin{table}[pos=htbp]
  \centering
  \caption{Gate outcomes as executed. Origin marks whether the criterion is in the frozen
  preregistration or was introduced by protocol/amendment. Three gate criteria failed (parser,
  severity $s2/s3$, recovery); the parser criterion's provenance is A-011, and the original G1
  criterion failed but the main run proceeded under the post-freeze amendment A-004, so it is
  disclosed separately rather than counted. The Weibull
  rejection is part of the G4 recovery failure, not a separate count. Recovery crude rates are
  run-risk-set Events/n (Qwen 72/264, Ministral 141/406, Llama 112/412),
  recomputed in \texttt{artifacts/analysis/g4\_recovery.json} (\texttt{transient\_s3} groups);
  the semantic risk-set re-estimation (Table~\ref{tab:recovery},
  \texttt{artifacts/analysis/g4\_recovery\_v2.json}) confirms the failure with wider margin.}
  \label{tab:gate-status}
  \footnotesize
  \setlength{\tabcolsep}{4pt}
  \begin{tabular}{>{\raggedright\arraybackslash}p{1.4cm}>{\raggedright\arraybackslash}p{6.4cm}>{\raggedright\arraybackslash}p{4.6cm}c}
    \toprule
    \textbf{Gate} & \textbf{Criterion (origin)} & \textbf{Observed} & \textbf{Verdict} \\
    \midrule
    G0 & $\ge880$ eligible; replay $\ge98\%$; 40-item human audit $\ge36/40$ (preregistered) & 5{,}516 eligible; 99.91\% replay; AI pre-audit 40/40 adopted, human verification deferred (D-021; A-005) & provisional \\
    G1 orig.\ & V-vs-$s0$ jitter $\le20\%$ (preregistered) & 20.3 / 36.7 / 45.6\% & \textbf{fail} \\
    G1 amend.\ & $s0$-vs-$s1$ $\le20\%$; parse $\ge95\%$; $s3$ 10--80\%; refusal $<5\%$ (A-004: post-freeze, not blind, D-041) & 10.5--13.3\%; parse 96.7--100\%; $s3$ 35.6--56.1\%; refusal 0--1.7\% & pass \\
    G3 slope & pooled $\beta>0$; $\ge2/3$ models same sign (preregistered) & $\hat\beta=0.877$; 3/3 positive & pass$^{\ast}$ \\
    G3 severity & $\alpha\ge.75$; $s2/s3$ $\ge90\%$ (preregistered) & $\alpha=.911$; 88.3\% (53/60) & \textbf{fail} \\
    G3 outcome & $\kappa\ge.80$; PPV $\ge.90$ (preregistered) & $\kappa=.987$; PPV 1.00, 9/9 (95\% lower bound .66) & pass \\
    G3 parser & failure $<1\%$ (protocol, not frozen prereg; A-011) & 205/13{,}200 (1.55\%) & \textbf{fail} \\
    G3 recovery agr.\ & $\kappa\ge.75$ (preregistered) & .920 / .917 & pass \\
    G4 & $\ge50\%$ of $s3$ risk set recovers by $H=4$; Weibull $p\ge.01$ (preregistered) & crude 27.3/34.7/27.2\% (72/264; 141/406; 112/412); KM 32.4/47.6/28.2\%; Weibull rejected in 2/3 primary groups & \textbf{fail} \\
    \bottomrule
  \end{tabular}
\end{table}
\noindent $^{\ast}$The slope passes as a measured shift among movers; because three gate criteria
failed, no causal or content-sensitivity claim is made.

\paragraph{G0: harvest and freeze}.\
G0 is CPU-only and costs \$0. It has five checks: supply of at least 880 eligible prefixes;
domain balance of at least 440 per domain before audit and 400 after audit; replay of at least
98\%; a 40-item human audit at 36/40 or better; and freeze integrity, meaning that message hashes,
observation hashes, perturbation records, the renderer version, the field-registry hash, and the
preregistration hash are recorded in \texttt{data/manifests/toolresult800\_v1.json}.
The implementation status as of 2026-11-15 is as follows. Supply passes with 5{,}516 eligible.
Replay passes with 5{,}511/5{,}516 (99.91\%). The cap-10 pool has 442 airline and 1{,}144 retail
prefixes, so the per-domain floor is met with a two-prefix airline margin. Rendering produced
4{,}400 arms and 2{,}640 perturbation records with 100\% $s=0$ value-set equality, and the
calibration set of 60 disjoint prefixes is selected. The freeze dry-run validates the manifest
shape: 800 plus 80 prefixes, 4{,}400 arms, 2{,}640 perturbations, and 60 calibration prefixes.
The freeze is complete: the manifest records 800 plus 80 prefixes, 4,400 arms, 2,640 perturbations,
and 60 calibration prefixes, with message hashes, observation hashes, renderer version,
field-registry hash, and preregistration hash. The 40-item sheet remains provisionally filled from
the AI-assisted mechanical pre-audit that reports 40/40 (D-021), with the author's final
verification deferred; the preregistration requires a human audit at 36/40, so that check is
reported as provisional rather than passed.

\paragraph{G1: calibration}.\

G1 is a harness check, not a result. It uses 60 prefixes (30 per domain), five arms, and three
models, for 900 generations, and the calibration prefixes are disjoint from the main 800 and stay
out of the primary analysis. The per-model thresholds are: parse rate at least 95\%; placebo
jitter ($s=0$ versus V) at most 20\%; $s=3$ switch rate between 10\% and 80\%; and format refusal
below 5\%. The gate also inspects the parser failure taxonomy, prompt truncation, context length,
latency, and whether $s=3$ is a ceiling. A model that refuses the forced next-decision format is
fixed with a prompt or grammar change; if it still refuses, it is excluded and the substitution
policy applies under a prereg amendment recorded before calibration results are used. No
main-run spend occurs before G1 passes. G1 passed on all three models only under the amended
negative-control criterion (A-004; $s0$-vs-$s1$ jitter 10.5--13.3\%); the original preregistered
V-vs-$s0$ format-jitter criterion failed at 20.3/36.7/45.6\% against the 20\% threshold and is
reported as a failed criterion in Table~\ref{tab:gate-status} and separately in
Table~\ref{tab:g1}. Amendment A-004 is post-freeze and not blind to the G1 calibration values
(D-041); the main run proceeded under A-004.

\paragraph{G2 to G4}.\

G2 is an operational stop rather than a scientific test. The runner checks cumulative GPU-hours
and storage after every shard and every 40 prefixes; on reaching 45 GPU-hours or 200~GB it stops
cleanly, flushes the ledger, and writes a stop record. Completed shards are analysed and missing
cells are reported, not imputed. G3 is the analysis gate on the slope and the agreement
statistics, and G4 is the recovery gate. Both are defined in Table~\ref{tab:gates} and both have
run: G3 passes on the slope but fails on the parser-failure and severity $s2/s3$ criteria, and G4
fails on recovery (Section~\ref{sec:results}). Algorithm~\ref{alg:matrix} summarizes the frozen
matrix construction.

\begin{algorithm}[htbp]
  \caption{Frozen fixed-prefix matrix and labels (G0--G2).}
  \label{alg:matrix}
  \begin{algorithmic}[1]
    \Require pinned $\tau$-bench trajectories $\mathcal{T}$, task databases $\mathcal{D}$,
    budgets $N=800$ and $N_{r}=80$, seed
    \State parse $\mathcal{T}$ into candidate post-tool decision points
    \State scrub direct identifiers; replay each candidate against $\mathcal{D}$
    \State retain candidates with exact replay and a value-used observation
    \State stratified sample $N+N_{r}$ by domain and donor; freeze $x_p$ and $o_p$
    \For{each frozen prefix $p$}
      \State $o_p^{(0)}\gets\Pi_0(o_p)$; verify value-set equality
      \For{$s\in\{1,2,3\}$}
        \State $o_p^{(s)}\gets\Pi_s(o_p)$; record changed paths, true/false values, and severity
      \EndFor
      \State render $R_a(o_p)$ for $a\in\{V,0,1,2,3\}$; record prompt hashes
    \EndFor
    \State freeze the manifest with message, observation, renderer, and preregistration hashes
    \State generate $d_{p,m,a}$ greedily for every $(p,m,a)$; write one ledger row per generation
    \State parse each row and assign $Y,Y^{A},Y^{F},Y^{U},Y^{O}$ by
    Equations~\ref{eq:switch}--\ref{eq:partition}
  \end{algorithmic}
\end{algorithm}



\paragraph{Environment}.\

All confirmatory inference runs on the approved university H200 MIG 1g.18gb slice: 18~GB of
VRAM, driver 580.178.04, CUDA 13.5, vLLM 0.33.0, and torch 2.14.0, in a project-local
environment. No system package is modified and no second GPU is used. The slice is smaller than
the 24~GB assumption in the original plan, which is why the fit check in
Section~\ref{sec:fit} precedes G1. The local RTX 3060 runs the CPU-side harvest, replay, and
analysis code and the smoke tests; it does not run study inference. Storage is expected to be
about 60~GB for three AWQ weights, caches, and traces, inside the 200~GB project allocation.

\paragraph{Models and fit checks}.\ \label{sec:fit}

Each model is loaded alone at 0.85 GPU memory utilization with a 24k context; the non-study
fit checks on synthetic prompts pass at batch 4 for all three (Table~\ref{tab:models}). In the run, the runner submits 8
prefixes per step, and the engine admits up to 8 concurrent sequences for Qwen3.5-9B and 4 for
Ministral-3-8B and Llama-3.1-8B, which ran out of memory at 8 concurrent sequences during
calibration. The per-model seq caps are the served configuration (D-016;
\texttt{configs/models.yaml}). A fit failure is handled by reducing batch or context first, then
by reducing $N$ or dropping the persistent arm under a logged decision; it is never handled by
booking a paid rental. Non-study fit checks are folded into Table~\ref{tab:models}.
\subsection{Evaluation Metrics}
\label{sec:metrics}

\paragraph{Agreement and audit statistics}.\

Agreement is reported with Krippendorff's $\alpha$ for the ordinal severity rating
\cite{krippendorff,hayes2007} and Cohen's $\kappa$ for the categorical outcome label
\cite{cohen1960,landis1977}.
For two raters with observed agreement $p_o$ and chance agreement $p_e$,
\begin{equation}
  \kappa=\frac{p_o-p_e}{1-p_e},
  \qquad
  p_o=\frac{1}{n}\sum_{i=1}^{n}\mathbf{1}[a_i=b_i],
  \qquad
  p_e=\sum_{c}\hat\pi_a(c)\,\hat\pi_b(c),
  \label{eq:kappa}
\end{equation}
where $\hat\pi_a(c)$ and $\hat\pi_b(c)$ are the two raters' category marginals. Krippendorff's
$\alpha$ generalizes the same idea to any number of raters and to ordinal codes,
\begin{equation}
  \alpha=1-\frac{D_o}{D_e},
  \qquad
  \delta^{2}(c,c')=
  \begin{cases}
    \mathbf{1}[c\neq c'], & \text{nominal},\\
    (c-c')^{2}, & \text{interval},
  \end{cases}
  \label{eq:alpha}
\end{equation}
where $D_o$ and $D_e$ are the observed and expected disagreement computed from the coincidence
matrix \cite{krippendorff}. Silent-propagation positive predictive value and recall are
\begin{equation}
  \mathrm{PPV}=\frac{TP}{TP+FP},
  \qquad
  \mathrm{Recall}=\frac{TP}{TP+FN},
  \label{eq:ppv}
\end{equation}
where $TP$ the automatic silent labels that the blind audit confirms, $FP$ the automatic labels it
rejects, and $FN$ the panel-confirmed silent items the coder misses. Switch-rate intervals use the
Wilson score interval,
\begin{equation}
  \frac{\hat p+\frac{z^{2}}{2n}\pm
        z\sqrt{\frac{\hat p(1-\hat p)}{n}+\frac{z^{2}}{4n^{2}}}}
       {1+\frac{z^{2}}{n}},
  \label{eq:wilson}
\end{equation}
where $z=1.96$ \cite{wilson1927}. Seed-repeat stability is the share of repeated generations whose decision signature is
unchanged. These statistics gate the claim rather than decorate it: a low $\kappa$ downgrades the
study to action-level $Y^{A}$, and a low $\alpha$ invalidates the dose interpretation even if the
slope is positive.

The minimum detectable effect is
\begin{equation}
  \mathrm{MDE}_{\beta}=
  \left(z_{1-\alpha/2}+z_{0.8}\right)\widehat{\mathrm{SE}}(\hat\beta),
  \qquad
  \mathrm{MDE}_{\mathrm{OR}}=e^{\mathrm{MDE}_{\beta}},
  \label{eq:mde}
\end{equation}
where $\alpha=.05$ is the two-sided level and power is $80\%$.

\paragraph{Generation envelope}.\

The main run crosses 880 frozen prefixes (800 primary analyzed, 80 reserve excluded)
$\times$ 5 arms $\times$ 3 models for the base matrix,
plus recovery and jitter. Table~\ref{tab:envelope} reconciles the planned envelope with the
executed rows. The base matrix yields
13{,}200 decisions (12{,}000 analyzed). Recovery continuation is conditional on switching at
the injected step, so the planned 19{,}200 transient rows were an upper bound; the executed
transient total is 6{,}430 rows (2{,}952/1{,}980/1{,}498 by model) plus 781 persistent rows
(339/239/203). The persistent regime was designed to corrupt the same field
on 25\% of prefixes (200) at $s=3$; in the executed runs the re-application found the corrupted
path in only a handful of steps, so the persistent prompts are byte-identical to transient for
97--100\% of shared steps and those groups are descriptive only. Jitter repeats use 80 prefixes
across all five arms to estimate decoding and prompt-format noise; they do not add power to the
primary estimate. The executed total is 21{,}611 rows at about 8.9 GPU-hours of the 45-hour
allocation (D-034 reconstruction, upper bound $\approx$8.95~h; \texttt{docs/DECISIONS.md}).

\begin{table}[pos=htbp]
  \centering
  \caption{Planned vs executed generation envelope. Recovery continuation is conditional on
  switching at the injected step, so planned recovery rows were an upper bound. GPU total is
  the D-034 stage-span reconstruction ($\approx$8.9~h of the 45-hour cap).}
  \label{tab:envelope}
  \small
  \begin{tabular}{>{\raggedright\arraybackslash}p{3.6cm}>{\raggedright\arraybackslash}p{4.6cm}rr}
    \toprule
    \textbf{Component} & \textbf{Planned calculation} & \textbf{Planned} & \textbf{Executed} \\
    \midrule
    Base decisions            & $800 \times 5 \times 3$ (run: $880 \times 5 \times 3$) & 12{,}000 & 13{,}200 \\
    Recovery, transient       & $800 \times 2 \times 4 \times 3$ (at-risk only)       & 19{,}200 & 6{,}430 \\
    Recovery, persistent      & $200 \times 1 \times 4 \times 3$ (at-risk only)       & 2{,}400  & 781 \\
    Jitter repeats            & $80 \times 5 \times 3$                                & 1{,}200  & 1{,}200 \\
    \midrule
    Total                     &                                                       & 34{,}800 & 21{,}611 \\
    GPU-hours                 & stop at 45 GPU-hours                                & ---       & $\approx$8.9 \\
    \bottomrule
  \end{tabular}
\end{table}

\paragraph{Execution, ledger, and stop rules}.\

Execution is sharded by model and domain. Each shard checkpoints every 40 prefixes and appends
one ledger row per generation: run ID, model, prefix, arm, step, prompt hash, output, parse
status, latency, GPU-hours, storage, and UTC timestamp. The runner is resumable after an H200
session interruption and never overwrites a completed row. A parser-error rate above 5\% in a
shard fails fast. The calibration ledger is retained and reported alongside the main ledger. The
run stops at 45 GPU-hours or 200~GB storage, whichever comes first; if the stop lands
mid-envelope, the analysis uses completed shards and the paper discloses the shortfall. No
imputation is used for missing arms, and the missingness rate is reported.

\paragraph{Analysis software and reproducibility}.\

The analysis code is versioned; run ledgers carry the manifest, prompt, and renderer
versions, and every analysis artifact records the producing commit. The primary conditional-logit fit, the prefix-cluster bootstrap, and the
recovery estimators are implemented in pure Python and tested on synthetic fixtures, so the
analysis can be reproduced without a GPU. A one-command validation target
(\texttt{make validate-release}, \texttt{scripts/18\_build\_release.py --validate}) reruns the five
CPU-only analysis producers, diffs the rebuilt JSONs against the released copies, rechecks every
manifest hash, reconciles the table-to-artifact map (\texttt{TABLE\_MAP.md}), and rescans the
package for identifiers. Manifests are the reproducibility spine: every table
in the paper traces to a manifest row, a ledger extract, or an analysis artifact via
\texttt{TABLE\_MAP.md}, and the release packager
refuses to build without a passed audit and a clean identifier re-scan. The release record binds
the preregistration, configuration, and renderer hashes to every bundled file. The workflow follows reproducibility
guidance that combines preregistration, artifact manifests, and leakage checks
\cite{nosek2018,munafo2017,pineau,kapoor2023}. The preregistration freeze (commit 1673856,
tag prereg-freeze-v1, hash \texttt{e0d107c6}; no external registry) and the deposit record are in
\texttt{prereg/DEPOSIT\_NOTE.md} \cite{toolresult800prereg}; amendments A-001--A-011 are
in \texttt{prereg/AMENDMENTS.md} (A-007--A-011 entered 2026-11-22 for the post-freeze analysis
record; A-010 is exploratory).

\section{Results}\label{sec:results}

The confirmatory analysis of the base matrix is complete. The run generated the full 880-prefix
matrix (13,200 base generations, five arms, three models, greedy decoding); the confirmatory
analysis uses the 800 primary prefixes and excludes the 80 reserve prefixes, as pre-registered
(A-007).
The recovery sequences and the 10\% seed-repeat shard are collected, and the blind audits are
complete in a label-free round; a multi-role agent panel, author-overseen, re-derived the ratings (D-043). This
section reports the pilot, the G0 harvest counts, the G1 calibration, the
primary estimate, silent propagation, recovery, and sensitivity. Results are reported as
measured; the three failed gate criteria (parser failure, severity $s2/s3$ decision share, and
recovery) are stated as failures.

\paragraph{Pilot evidence (screening only)}.\ \label{sec:pilot}

The pilot used $N=60$ synthetic WebArena-lite / $\tau$-style prefixes, \texttt{qwen3.5:4b}, greedy
decoding, and three collapsed action codes (trust, re-call, flag). It compared a truthful arm, a
placebo arm, and an error arm. Table~\ref{tab:pilot} reports the recorded results. The consult
rate of 90\% shows that the model reads the observation. The placebo jitter of zero shows that
the format change alone was not associated with a different collapsed action code in this sample. The error-switch
rate of 8.3\% shows that error content was associated with a different next action. The paired McNemar test is the
design-matched comparison and is marginal at $p=.0625$; the recorded 95\% interval is an unpaired
Wald interval and is compatible with a 1.3 percentage-point effect. A 12-item human spot-check
agreed with all 12 labels, but the letter coding missed some silent fact propagation, so 8.3
percentage points is an action-level lower bound rather than a content-sensitivity estimate.
The pilot is not evidence for the confirmatory claim. It used one 4B model, synthetic prefixes,
coarse action coding, and no recovery data, and it cannot distinguish a graded law from a handful
of flips. Its role is to show that the manipulation is runnable and that the confirmatory design
is worth the GPU time.

\begin{table}[pos=htbp]
  \centering
  \caption{Pilot results on 60 synthetic prefixes; excluded from confirmatory analysis.}
  \label{tab:pilot}
  \small
  \begin{tabular}{>{\raggedright\arraybackslash}p{7.4cm}r}
    \toprule
    \textbf{Quantity} & \textbf{Result} \\
    \midrule
    Consult rate                              & 54/60 = 90.0\% \\
    Placebo jitter (placebo vs truthful)      & 0/60 = 0\% \\
    Error switch (error vs truthful)          & 5/60 = 8.3\% \\
    Gap, 95\% CI                              & 8.3 pp [1.3, 15.3] \\
    McNemar                                   & $b=5$, $c=0$; exact $p=.0625$ \\
    Human spot-check                          & 12/12 labels agree \\
    \bottomrule
  \end{tabular}
\end{table}

\paragraph{G0 harvest and freeze status}.\

The G0 implementation counts are in Table~\ref{tab:supply}. The 6{,}554 planning figure is
superseded by the stricter operational value-used rule. Supply, domain balance, replay, and
rendering meet their thresholds. The manifest is frozen at
\texttt{data/manifests/toolresult800\_v1.json} (SHA-256 sidecar
\texttt{toolresult800\_v1.sha256.json}); the frozen prefix payload is
\texttt{data/frozen/prefixes.jsonl}. The 40-item audit is provisionally filled from an AI-assisted
mechanical pre-audit at 40/40 under decision D-021, with the authors' final verification deferred;
the preregistration requires a human audit at 36/40 or better. The audit deviation is reported
here rather than in a footnote, because a reader should be able to see exactly how the gate was
filled. Final author verification is deferred under D-021; any replacement rows and the
verification date will be logged in the manifest record.

\paragraph{G1 calibration}.\

G1 ran under the fixed harness: \texttt{think} was excluded, decoding was constrained to a
valid tool-or-final-answer JSON object, the decision budget was 1{,}024 tokens, and each model ran
in a separate process. All three models passed. Table~\ref{tab:g1} reports the per-model values.
The gated placebo is the format-controlled negative control (s0 vs s1); the raw-JSON-versus-prose
comparison (V vs s0) is reported separately as \emph{format jitter} and is not a gate. The format
effect is large (20.3--45.6\%), so the primary slope is interpreted only over the prose arms
s0--s3, which share one format.

\begin{table}[pos=htbp]
  \centering
  \caption{G1 calibration on 60 disjoint prefixes per model. The gated placebo is the
  format-controlled $s0$-vs-$s1$ negative control (A-004, post-freeze and not blind per D-041);
  the V-vs-$s0$ column is the failed original criterion, shown for format sensitivity and not gated.}
  \label{tab:g1}
  \small
  \begin{tabular}{lccccc}
    \toprule
    \textbf{Model} & \textbf{Parse} & \textbf{Neg.\ control} & \textbf{Format (V/s0)} & \textbf{s=3 switch} & \textbf{Refusal} \\
    \midrule
    Qwen3.5-9B     & 98.3\% (295/300) & 11.9\% & 20.3\% & 35.6\% & 1.7\% \\
    Ministral-3-8B & 100\% (300/300)  & 13.3\% & 36.7\% & 55.0\% & 0.0\% \\
    Llama-3.1-8B   & 96.7\% (290/300) & 10.5\% & 45.6\% & 56.1\% & 0.0\% \\
    \bottomrule
  \end{tabular}
\end{table}

\subsection{Primary severity--switch estimate}

Table~\ref{tab:hypotheses} maps each preregistered hypothesis to its observed value and status;
the slope and tier paragraphs below give the H1--H3, H5--H6 detail.

\begin{table}[pos=htbp]
  \centering
  \caption{Hypothesis status (preregistration \S2). H1--H5 are confirmatory, H6--H7 exploratory.
  H5 is reported here ($s=3$ vs $s=1$ paired exact McNemar: pooled $b=431$, $c=93$, two-sided
  $p=5.3\times10^{-53}$; per-model two-sided $p=1.98\times10^{-9}$ (Llama-3.1-8B),
  $9.0\times10^{-22}$ (Ministral-3-8B), $7.5\times10^{-27}$ (Qwen3.5-9B);
  \texttt{artifacts/analysis/prereg\_secondary.json}). The slope rows pass as a measured
  shift among movers; no causal or content-sensitivity claim is made because three gate criteria
  failed (Table~\ref{tab:gate-status}). H4 crude counts are Events/n (72/264; 141/406; 112/412)
  from \texttt{artifacts/analysis/g4\_recovery.json} (\texttt{transient\_s3} groups).}
  \label{tab:hypotheses}
  \small
  \begin{tabular}{>{\raggedright\arraybackslash}p{1.1cm}>{\raggedright\arraybackslash}p{6.8cm}>{\raggedright\arraybackslash}p{4.6cm}c}
    \toprule
    \textbf{ID} & \textbf{Hypothesis} & \textbf{Observed} & \textbf{Status} \\
    \midrule
    H1 & Pooled slope $\beta>0$ over $s=0,\dots,3$ & $\hat\beta=0.877$ [0.770, 0.998], $p=2.3\times10^{-49}$ & supported$^{\ast}$ \\
    H2 & Per-model slope positive in $\ge2/3$ models, same sign & 3/3 positive, Holm-significant & pass \\
    H3 & Silent propagation a non-zero share at $s=2,3$ & 419/1{,}964 = 21.3\% of critical-arm switches & supported$^{\dagger}$ \\
    H4 & $\ge50\%$ of $s3$ risk set recovers by $H=4$; Weibull $p\ge.01$ & crude 27.3/34.7/27.2\%; Weibull rejected 2/3 & \textbf{fail} \\
    H5 & $s=3$ switch rate exceeds $s=1$ & McNemar $b=431$, $c=93$, $p=5.3\times10^{-53}$ & supported \\
    H6 & Donor/task slopes directionally consistent & donor interaction $p=.61$; domain 0.74/0.99, interaction $p=.035$ (exploratory) & consistent \\
    H7 & Model $\times$ severity, task $\times$ severity not claimed & reported exploratory only & --- \\
    \bottomrule
  \end{tabular}
\end{table}
\noindent $^{\ast}$As a measurement estimate only (gate fail action). $^{\dagger}$Descriptive
secondary analysis, gated on the outcome audit (pass); silent share is a lower bound
(observed 21.3\%, recall .75, recall-adjusted $\approx$28.5\%).

\begin{table}[pos=htbp]
  \centering
  \caption{Primary severity-switch estimates by model and pooled. Per-model intervals are
  Wald 95\% CIs from the prefix-cluster SE; the pooled interval is the prefix-cluster
  bootstrap ($2{,}000$ replicates). Figure~\ref{fig:forest} plots the Wald intervals with
  the pooled bootstrap band.}
  \label{tab:primary}
  \small
  \begin{tabular}{lrrrrr}
    \toprule
    \textbf{Model} & $\hat\beta$ & SE (cluster) & $z$ & $p$ & \textbf{95\% CI} \\
    \midrule
    Qwen3.5-9B     & 1.079 & 0.114 & 9.45 & $3.4\times10^{-21}$ & [0.86, 1.30] \\
    Ministral-3-8B & 0.915 & 0.090 & 10.12 & $4.7\times10^{-24}$ & [0.74, 1.09] \\
    Llama-3.1-8B   & 0.649 & 0.100 & 6.46 & $1.0\times10^{-10}$ & [0.45, 0.85] \\
    \midrule
    Pooled         & 0.877 & 0.059 & 14.77 & $2.3\times10^{-49}$ & [0.770, 0.998]$^{\ast}$ \\
    \bottomrule
  \end{tabular}
\end{table}
\noindent $^{\ast}$Bootstrap interval; the pooled Wald interval is [0.76, 0.99].

\begin{figure}[pos=htbp]
  \centering
  \figfile[0.95\linewidth]{figures/fig5.png}
  \caption{Per-model and pooled severity slopes. Error bars are Wald 95\% CIs from the
  prefix-cluster SEs; the shaded band and diamond show the pooled prefix-cluster bootstrap
  interval [0.770, 0.998] reported in Table~\ref{tab:primary}.}
  \label{fig:forest}
\end{figure}

Table~\ref{tab:primary} reports the primary estimate, the pooled conditional-logistic slope of the switch indicator of
Equation~\ref{eq:switch} on severity, with the per-model slopes as the replication gate
(Figure~\ref{fig:forest}). The
base matrix is complete (13{,}200 generations, 880 prefixes, five arms, three models, greedy
decoding), so this is the pre-registered estimate, not a partial one. The pooled slope is $\hat\beta = 0.877$ with a prefix-cluster standard error of $0.059$
($z = 14.77$, $p = 2.3\times10^{-49}$). The prefix-cluster bootstrap interval is
$[0.770, 0.998]$ ($2{,}000$ replicates, $p = 0.001$ two-sided). All three per-model slopes
are positive and Holm-significant, so the pre-registered replication gate is met. The slope is
the log odds of a decision switch per severity level, conditional on the prefix, so an odds
ratio of $e^{0.877} \approx 2.40$ per level, but that summary hides the shape: the tier odds
ratios against $s0$ are 1.27 (s1), 4.60 (s2), and 11.21 (s3) (Table~\ref{tab:tiers} and
Figure~\ref{fig:tiers}), and the pre-registered linearity diagnostics reject a constant gradient. The estimate uses 617 prefix$\times$model strata
(423 unique prefixes) and 2{,}458 slope rows from the 800 primary prefixes; the 80 reserve
prefixes are excluded (A-007). Of the 9{,}075 labelled prose-arm cells, 2{,}297 form prefix$\times$model
strata and 617 vary across arms; the conditional logit drops the 1{,}680 concordant strata, which
contribute no information, so the estimate is a within-stratum effect identified from prefixes
whose decision varies across arms.
The switch typology is reported separately (\texttt{artifacts/analysis/switch\_typology.json}):
47.8\% of strata never switch, 23.9\% switch in every arm, 19.0\% are dose-responsive
($s0=s1=0$, $s2$ or $s3=1$), and 9.3\% are placebo switchers (Table~\ref{tab:typology}). The claim is bounded to the prose
arms s0--s3: the raw-JSON arm V is excluded because the format effect is large (20.3--45.6\%,
D-026). An earlier analysis under the pre-fix automatic coder gave $\hat\beta = 0.684$
[0.584, 0.797] (artifact not retained; A-006); the same post-fix coder on all 880 prefixes gives 0.867 [0.760, 0.988]
(\texttt{g3\_analysis\_all880.json}), so the coder change accounts for most of the movement and
the reserve exclusion for the remainder. The blind audit then found and corrected the labeling
defect (D-028, A-006), and the estimate above is the corrected one. All three are positive and all
are reported. The G3 gate did not pass: the protocol parser-failure criterion is 1.55\% (205/13{,}200; 3.45\% including
refusals) against the 1\% threshold (a protocol criterion from \texttt{docs/GATES.md}, absent from
the frozen prereg G3 row; A-011). Under the gate's fail action the paper reports a measured shift among
movers and makes no causal or content-sensitivity claim; the slope is not a confirmed effect. The recovery endpoint is reported separately; the audits gate the label-based
decomposition, not the switch indicator. Estimates trace to
\texttt{artifacts/analysis/g3\_analysis.json}.

\begin{table}[pos=htbp]
  \centering
  \caption{Tier odds ratios against the prose placebo.}
  \label{tab:tiers}
  \footnotesize
  \setlength{\tabcolsep}{4pt}
  \begin{tabular}{lccc r}
    \toprule
    \textbf{Sample} & \textbf{OR $s=1$ [95\% CI]} & \textbf{OR $s=2$ [95\% CI]} & \textbf{OR $s=3$ [95\% CI]} & \textbf{$n$ groups/rows} \\
    \midrule
    Qwen3.5-9B     & 1.01 [0.65, 1.56] & 6.49 [4.25, 10.90] & 17.12 [10.10, 33.00] & 201/799 \\
    Ministral-3-8B & 1.39 [1.00, 1.99] & 4.98 [3.43, 7.34]  & 12.77 [8.18, 20.80]  & 246/979 \\
    Llama-3.1-8B   & 1.36 [1.03, 1.80] & 3.12 [2.10, 4.96]  & 6.42 [4.04, 11.65]   & 170/680 \\
    \midrule
    Pooled         & 1.27 [1.04, 1.55] & 4.60 [3.58, 5.99]  & 11.21 [8.38, 15.30]  & 617/2{,}458 \\
    \bottomrule
  \end{tabular}
\end{table}

\begin{figure}[pos=htbp]
  \centering
  \figfile[0.98\linewidth]{figures/fig6.png}
  \caption{Tier odds ratios against $s=0$ on a log scale.}
  \label{fig:tiers}
\end{figure}

\paragraph{Pre-registered diagnostics}.\

The categorical-vs-linear likelihood-ratio test (Table~\ref{tab:diagnostics}) rejects linearity pooled
($\chi^2(2) = 25.03$, $p = 3.7\times10^{-6}$; Qwen $p = 1.0\times10^{-4}$, Ministral $p=.017$,
Llama $p=.274$), and the quadratic term is positive ($b = 0.143$, cluster SE $0.036$,
$p = 7.4\times10^{-5}$). Tier odds ratios against $s0$ are 1.27 [1.04, 1.55] (s1), 4.60
[3.58, 5.99] (s2), and 11.21 [8.38, 15.30] (s3). The $s2$-vs-$s3$ McNemar contrast is
primary per model ($p=1.1\times10^{-7}$, $2.6\times10^{-10}$,
$5.8\times10^{-5}$) and secondary pooled ($p = 2.1\times10^{-20}$ over 2{,}251
model-prefix pairs); the pooled value survives a prefix-cluster design-effect grid to
$\rho=0.5$ and a model-stratified check (Section~\ref{sec:method};
\texttt{artifacts/analysis/inference\_robustness.json}). No new bootstrap LRT was run
(row-level refits need the run ledgers, which are not in the repo); the resampling
robustness is the 2{,}000-replicate prefix-cluster bootstrap tier-OR intervals, whose
$s2$/$s3$ intervals exclude one in every model while the Qwen $s1$ interval includes one. The pattern is a threshold at $s2$, not a graded per-level law. Domain
slopes are 0.74 [0.59, 0.93] (airline) and 0.99 [0.83, 1.17] (retail); the domain interaction is
$p=.035$ (exploratory, driven by Ministral) and the donor interaction is null ($p=.61$). Under
lenient tool-or-verb coding the pooled slope is 0.43 [0.30, 0.57] versus 0.88 [0.77, 1.00] under
the strict signature: coding changes the magnitude, not the sign. Excluding the placebo,
$\hat\beta = 1.03$ [0.88, 1.19]. At the analyzed sample the realized pooled MDE is 0.166
logit/level (OR 1.18, 80\% power), and realized power at the observed effect is 1.0.
Diagnostics trace to \texttt{artifacts/analysis/prereg\_secondary.json}, with the $s1$ tier
near one and the jump at $s2$.
The $s1$ negative control is not inert and is reported as a low dose with contamination
rather than a clean zero: the conditional $s1$-vs-$s0$ odds ratio is 1.27 [1.04, 1.55]
($p=.019$), while the marginal $s1$--$s0$ contrast is $+1.1$ points (Newcombe 95\%
$[-1.6, +3.8]$), not significant on the margin. The audit rates 3/30 $s1$ items
decision-changing (majority; Wilson 95\% [3.5, 25.6]\%), so the non-critical arm carries
realized critical content at the margin. The threshold reading is therefore a gradient from
$s0$ through a contaminated $s1$ to $s2$, steepening at $s2$, not a step from a true zero.

\begin{table}[pos=htbp]
  \centering
  \caption{Linearity and paired diagnostics by sample. Per-model McNemar tests are the
  primary paired contrast; the pooled McNemar $p$ is secondary (prefix-cluster caveat in
  Section~\ref{sec:method}).}
  \label{tab:diagnostics}
  \footnotesize
  \setlength{\tabcolsep}{4pt}
  \begin{tabular}{lccccc}
    \toprule
    \textbf{Sample} & \textbf{Cat.\ LRT $\chi^2(2)$} & \textbf{$p$} & \textbf{Quadratic $b$ (SE)} & \textbf{$p$} & \textbf{McNemar $b/c$ ($p$)} \\
    \midrule
    Qwen3.5-9B     & 18.29 & $1.0\times10^{-4}$ & 0.171 (0.067) & .011 & 70/20 ($1.1\times10^{-7}$) \\
    Ministral-3-8B & 8.18  & .017               & 0.135 (0.058) & .020 & 73/15 ($2.6\times10^{-10}$) \\
    Llama-3.1-8B   & 2.59  & .274               & 0.101 (0.052) & .054 & 47/15 ($5.8\times10^{-5}$) \\
    \midrule
    Pooled         & 25.03 & $3.7\times10^{-6}$ & 0.143 (0.036) & $7.4\times10^{-5}$ & 190/50 ($2.1\times10^{-20}$) \\
    \bottomrule
  \end{tabular}
\end{table}

\paragraph{Missingness and the population-average contrast}.\

Cell loss is 5.47\% (525/9{,}600 prose-arm cells): 326 unparsed prose-arm generations
(148/104/74 for Qwen/Ministral/Llama) plus 199 parsed prose-arm rows (85/67/47) dropped for a
missing parsed V reference (\texttt{artifacts/analysis/missingness\_marginal.json}). Of the
13{,}200 step-0 rows, 205 are hard parse errors (1.55\%) and 250 are refusals, so invalid
output including refusals is 455/13{,}200 (3.45\%). Per arm, 97.0--98.6\% of parsed rows
enter the analysis, and 2.3--4.9\% of prefixes lack a parsed V reference
(Table~\ref{tab:missing}; the column is a share of prefixes). The
conditional slope is estimated on 617 prefix$\times$model groups; 649 of the 2{,}297 strata
(28.3\%) vary across arms, while 47.8\% never switch and 23.9\% switch in every arm, including
the placebo. The
population-average paired contrasts over all labels are +10.2 points for $s2$ vs $s0$
[8.7, 11.7], +16.3 for $s3$ vs $s0$ [14.5, 18.0], and +6.2 for $s3$ vs $s2$ [4.9, 7.6]. The
conditional slope and the marginal contrast therefore agree on direction and on the threshold
shape; the slope should be read as a within-stratum (mover) effect, not a population-average
effect. Table~\ref{tab:armrates} gives the descriptive per-arm switch rates and
Figure~\ref{fig:dose} plots them with Wilson intervals. They rise
monotonically with severity for all three models, with $s0$ and $s1$ nearly flat, which is the
shape the slope summarizes; the s0 placebo rate is the format-jitter floor reported in
Table~\ref{tab:g1} (calibration 20.3--45.6\%; main-run $s0$ rates 19.8/31.9/38.1\%).
Parsed rows are 776--789 per arm (Table~\ref{tab:missing}); labeled rows are 740--770 per arm
after the V-unparsed filter
(\texttt{artifacts/analysis/g3\_analysis.json}).

\begin{table}[pos=htbp]
  \centering
  \caption{Analysis coverage by model and arm.}
  \label{tab:missing}
  \small
  \begin{tabular}{lccccr}
    \toprule
    \textbf{Model} & \textbf{s0} & \textbf{s1} & \textbf{s2} & \textbf{s3} & \textbf{V-unparsed (\%)} \\
    \midrule
    Qwen3.5-9B     & 97.4\% & 97.4\% & 97.1\% & 97.0\% & 4.9\% \\
    Ministral-3-8B & 98.0\% & 97.8\% & 97.8\% & 97.8\% & 3.4\% \\
    Llama-3.1-8B   & 98.6\% & 98.6\% & 98.5\% & 98.3\% & 2.3\% \\
    \bottomrule
  \end{tabular}
\end{table}

\begin{table}[pos=htbp]
  \centering
  \caption{Switch typology across prefix-by-model strata.}
  \label{tab:typology}
  \small
  \begin{tabular}{lrrrr}
    \toprule
    \textbf{Sample} & \textbf{Never} & \textbf{Always} & \textbf{Dose-responsive} & \textbf{Placebo switcher} \\
    \midrule
    Qwen3.5-9B     & 429 (56.7\%) & 111 (14.7\%) & 161 (21.3\%) & 56 (7.4\%) \\
    Ministral-3-8B & 320 (41.7\%) & 190 (24.7\%) & 168 (21.9\%) & 90 (11.7\%) \\
    Llama-3.1-8B   & 350 (45.3\%) & 248 (32.1\%) & 107 (13.9\%) & 67 (8.7\%) \\
    \midrule
    Pooled         & 1{,}099 (47.8\%) & 549 (23.9\%) & 436 (19.0\%) & 213 (9.3\%) \\
    \bottomrule
  \end{tabular}
\end{table}

\begin{table}[pos=htbp]
  \centering
  \caption{Descriptive switch rates by arm and model.}
  \label{tab:armrates}
  \small
  \begin{tabular}{lrrrr}
    \toprule
    \textbf{Model} & \textbf{s=0} & \textbf{s=1} & \textbf{s=2} & \textbf{s=3} \\
    \midrule
    Qwen3.5-9B     & 19.8\% & 19.9\% & 31.0\% & 38.4\% \\
    Ministral-3-8B & 31.9\% & 33.6\% & 44.4\% & 51.9\% \\
    Llama-3.1-8B   & 38.1\% & 39.7\% & 44.8\% & 49.0\% \\
    \bottomrule
  \end{tabular}
\end{table}

\begin{figure}[pos=htbp]
  \centering
  \figfile[0.95\linewidth]{figures/fig7.png}
  \caption{Switch rates by arm and model with 95\% confidence intervals.}
  \label{fig:dose}
\end{figure}

The claim is conditional on the sign pattern: a positive pooled slope with fewer than two
positive per-model slopes is a null result under the pre-registered decision rule, not a partial
success. Here all three are positive, so the slope stands. The severity, outcome, and recovery
audits were re-run label-free (D-039; A-008): severity $\alpha = .911$ (nominal; interval .959),
outcome $\kappa = .987$, silent-propagation PPV $= 1.00$ (9/9; exact 95\% lower bound .66) with
silent recall $.75$ (the coder misses three panel-confirmed silent items), and recovery
$\kappa = .920$ (step) and $.917$ (endpoint). The severity audit misses one of its five criteria: 88.3\% of
the design's $s2/s3$ items are rated decision-changing against the 90\% threshold (Table~\ref{tab:audit}; 7 majority-flagged items, 11--16 across verifier passes,
listed in \texttt{artifacts/audit/blind\_audit\_result.json}); the other severity criteria pass
(Jonckheere--Terpstra $p<10^{-4}$, $s1$ non-decision-changing 90.0\%, $s3$ absurd 3.3\%). The
earlier label-visible round ($\alpha = .919$, $\kappa = .975$, PPV $= 1.00$, recovery $.787$) is
archived as a sensitivity. An audit criterion did fail, and the claim is downgraded accordingly
rather than smoothed over.

\begin{table}[pos=htbp]
  \centering
  \caption{Pre-registered audit gates and observed values.}
  \label{tab:audit}
  \small
  \begin{tabular}{lrrl}
    \toprule
    \textbf{Criterion} & \textbf{Observed} & \textbf{Gate} & \textbf{Verdict} \\
    \midrule
    Severity Krippendorff $\alpha$ & .911 / .959 & $\ge .75$ & pass \\
    Severity Jonckheere--Terpstra & $J=5192$, $p<10^{-4}$ & $p<.01$ & pass \\
    $s=1$ non-decision share & 27/30 (90.0\%) & $\ge 90\%$ & pass \\
    $s=2/s=3$ decision share & 53/60 (88.3\%) & $\ge 90\%$ & \textbf{fail} \\
    $s=3$ absurd share & 1/30 (3.3\%) & $\le 20\%$ & pass \\
    Outcome Cohen's $\kappa$ & .987 & $\ge .80$ & pass \\
    Silent-propagation PPV & 9/9 (1.00; exact 95\% lower bound .66) & $\ge .90$ & pass \\
    Recovery Cohen's $\kappa$ & .920 / .917 & $\ge .75$ & pass \\
    Parser failure (G3) & 205/13{,}200 (1.55\%) & $<1\%$ & \textbf{fail} \\
    \bottomrule
  \end{tabular}
\end{table}

The miss is concentrated at $s2$: 25/30 $s2$ items are rated decision-changing (83.3\%) against
28/30 at $s3$ (93.3\%). The seven majority-flagged items are ones the donor continuation never
reads: a flight status, four route leaves, a cabin, two payment-type leaves, and insurance. The
range across passes is 7/60 (majority; 5 at $s2$, 2 at $s3$), 11/60 (sheet-visible verifier;
7 at $s2$, 4 at $s3$), and 16/60 (blind verifier; 11 at $s2$, 5 at $s3$), i.e.\ 12--27\% of
audited critical-arm stimuli are donor-unused (Wilson 95\% [5.8, 22.2]\%, [10.6, 29.9]\%,
[17.1, 39.0]\% respectively). In at least
three, the design's automatic usage test matched a short token in prose (``business'' in
``business days'', ``no'' in ``no additional charge'') or a route summary that restates an
unchanged top-level field. The same test flags three $s1$ items as decision-changing (each
verifier keeps one of the three), so the
noise runs both ways. A bounded dilution sensitivity
(\texttt{artifacts/analysis/severity\_sensitivity.json}) re-estimates the marginal tier odds
ratios treating donor-unused critical-arm rows as switching at the pooled $s1$ rate: the
observed marginal ORs of 1.56 ($s2$) and 2.02 ($s3$) rise to 1.64/2.19, 1.70/2.31, and
1.79/2.52 under $m=7/60$, $11/60$, and $16/60$, and the marginal $s1$--$s2$ gap widens from
9.0 to 12.3 points at the blind-verifier bound. Correcting the
assigned doses therefore steepens the $s1$--$s2$ contrast under every pass; the observed threshold is not an
artifact of this audit failure. The estimand is labelled assigned-severity throughout: the
conditional tier ORs 1.27/4.60/11.21 are contrasts over assigned arms, attenuated
(conservative) for realized criticality, and the per-level magnitudes are read as
descriptive rather than as a calibrated dose--response function.

An agent verification panel re-derived the ratings from the same packets (D-043; Table~\ref{tab:panel}). On outcome, the
verifier agreed with the blind majority on 158/160 items and confirmed the silent-propagation PPV
(9/9 automatic flags); it identified three further silent items the coder missed (OUT-029, OUT-078,
OUT-156), so silent recall is .75. On recovery, the verifier agreed with the majority on 66/80,
but all 14 differences were rater over-credits, and it reproduced the automatic endpoint and step
on 80/80 sequences. On severity, a sheet-visible pass agreed on 114/120 and a fully blind pass on
108/120; the two passes agree with each other on 112/120, and both place the $s2/s3$
decision-changing share below the gate (81.7\% and 73.3\%), with the disagreement concentrated at
the 1-vs-2 usage boundary (redundant route leaves, category words, unselected options). The
pre-registered gate remains the blind majority's 88.3\% and fails; the panel does not overturn it.
Two second raters independently re-rated a stratified 120-item subset from the same packets and
agreed at $\alpha = .917$ (severity), $\kappa = .942$ (outcome), and $\kappa = .918$ / $.907$
(recovery endpoint / step); against the primary severity consensus, $\alpha = .946$. The sessions
share one underlying model family, and two of the three severity sessions are identical on
116/120 items, so $\alpha = .911$ is consistency among near-duplicate readings, not independent
human reliability. The post-fix audit is explicitly in-sample: the corrected coder was conformed
on the same 160 outcome items later used to certify $\kappa=.987$ and PPV $=1.00$, so no fresh
held-out check exists. The held-out plan is a 60--120-item human audit on reserve prefixes in
future work; until then the agreement statistics are reproducibility evidence, not independent
validation. The panel record is
\texttt{artifacts/audit/verification/panel\_record.json}.

\begin{table}[pos=htbp]
  \centering
  \caption{Agent verification panel agreement and corrections.}
  \label{tab:panel}
  \small
  \begin{tabular}{lrrl}
    \toprule
    \textbf{Audit} & \textbf{Agreement} & \textbf{Corrections} & \textbf{Panel check} \\
    \midrule
    Outcome            & 158/160 & 2  & PPV 9/9; recall .75 \\
    Recovery           & 66/80   & 14 & automatic endpoint/step 80/80 \\
    Severity (visible) & 114/120 & 6  & $s2/s3$ decision share 81.7\% \\
    Severity (blind)   & 108/120 & 12 & $s2/s3$ decision share 73.3\% \\
    \bottomrule
  \end{tabular}
\end{table}

\paragraph{Silent propagation and the outcome decomposition}.\

The descriptive secondary analysis reports the share of decision changes that are silent
propagation, separately for $s=2$ and $s=3$, together with the detection and action-switch shares
and the exclusive typology of Equation~\ref{eq:partition}. Of the 1{,}964 decision switches on the
critical arms ($s=2$, $s=3$), 42.7\% are action switches, 1.5\% carry a flag without an action
switch, at least 21.3\% [19.6, 23.2] are silent propagation (the model repeats the injected false value without flagging
it; a lower bound because the coder misses panel-confirmed silent items --- recall .75 [.47, .91],
recall-adjusted share $\approx$28.5\%, range 23.4--45.6\% over the recall CI), and 34.4\% are other value-level changes not attributable to the injected false value
($Y^{O}$). The per-model silent share of switches is 24.4\% (Qwen3.5-9B), 23.3\%
(Ministral-3-8B), and 17.2\% (Llama-3.1-8B). Detection is rare in all three (0.1--1.9\% of
critical-arm rows), so silent propagation is the dominant \emph{identified} failure mode, not the
majority of switches; it is the quantity that end-to-end benchmarks cannot see.
The per-arm switch rates are in Table~\ref{tab:armrates}. This subsection is descriptive and is
not gated on a $p$-value, but it is gated on the outcome audit. Before the conformance fix the
automatic silent labels had a PPV of .76 (13/17); the audit found that the coder fired on outputs
identical to the reference and ignored the decision level, and after the fix (D-028) the PPV is
1.00 (all nine audit items confirmed by both raters). The audit changed 27 of 160 automatic
labels.

\paragraph{Which fields propagate silently}.\

An exploratory breakdown of the $s=2$ switches by perturbed field class
(Table~\ref{tab:fieldtypes}; \texttt{artifacts/\allowbreak analysis/\allowbreak
field\_taxonomy.json}; Figure~\ref{fig:fields} plots the four largest classes) shows that identifiers and money are the silent-propagation engines: 45.0\% and 45.6\% of their switches are
unflagged, action-preserving copies of the false value, versus 6.9\% (entity) and 6.6\% (status).
Entity and status fields instead flip the action in about half of their switches. The clearest
single path is the reservation identifier: 58.2\% of its $s=2$ perturbations are associated with a different decision and
51.9\% of those switches are silent. Flags are near-absent (0.8\% of $s=2$ rows). The three
smaller classes complete the partition: temporal 10/17 switches (58.8\%; 4/10 silent),
count 5/19 (26.3\%; 0/5 silent), and other 110/247 (44.5\%; 38/110 silent). At fixed
severity the field-class spread in switch rates is 32.5 percentage points (58.8\% temporal
vs 26.3\% count over all seven classes), while the criticality
contrast adds 10.1 points over a render floor of 30.1\%; field type is therefore at least as
informative as the critical/non-critical label. Wilson 95\% intervals are in the same artifact.

\begin{table}[pos=htbp]
  \centering
  \caption{Field-type breakdown of $s=2$ switches (all seven classes).}
  \label{tab:fieldtypes}
  \small
  \begin{tabular}{lrrr}
    \toprule
    \textbf{Field class} & $n$ & \textbf{Switch rate} & \textbf{Silent share of switches} \\
    \midrule
    Identifier & 388 & 53.9\% & 45.0\% \\
    Money      & 223 & 40.4\% & 45.6\% \\
    Entity     & 955 & 36.6\% & 6.9\% \\
    Status     & 418 & 32.8\% & 6.6\% \\
    Temporal   & 17  & 58.8\% & 40.0\% \\
    Count      & 19  & 26.3\% & 0.0\% \\
    Other      & 247 & 44.5\% & 34.5\% \\
    \midrule
    All $s=2$  & 2{,}267 & 40.2\% & 23.1\% \\
    \bottomrule
  \end{tabular}
\end{table}

\begin{figure}[pos=htbp]
  \centering
  \figfile[0.95\linewidth]{figures/fig8.png}
  \caption{Field-type switch rates and silent shares at $s=2$ (four largest classes;
  all seven in Table~\ref{tab:fieldtypes}).}
  \label{fig:fields}
\end{figure}

\subsection{Wrong-direction (utility) endpoint}

The switch indicator records deviation from the donor verbatim continuation, not task error.
The perturbation records supply the missing direction: every critical-arm ($s=2$, $s=3$)
record in the frozen manifest carries \texttt{expected\_action\_change=true} (1{,}760 records;
asserted fail-loud in \texttt{scripts/17\_wrong\_direction.py}), so a critical-arm switch that
is not an explicit flag moved in the anticipated direction class. The endpoint classifies each
of the 1{,}964 critical-arm switches exclusively: wrong-direction (action switch $C^{A}$ or
silent propagation $C^{U}$), caught (flag-only $C^{F}$), or ambiguous (other change $C^{O}$).
It is descriptive, with Wilson 95\% intervals, and traces to
\texttt{artifacts/analysis/wrong\_direction.json}.

Of the 1{,}964 critical-arm switches, 1{,}258 (64.1\% [61.9, 66.1]) are wrong-direction: 839
action switches plus the 419 confirmed silent propagations; 30 (1.5\% [1.1, 2.2]) are caught
by an explicit flag and 676 (34.4\%) are ambiguous (Table~\ref{tab:wrongdir}). The
wrong-direction share is directionally stable across severity (65.6\% [62.5, 68.7] at $s=2$,
62.7\% [59.7, 65.5] at $s=3$) and exceeds half in every model (Qwen3.5-9B 72.5\% [68.5, 76.2],
Ministral-3-8B 63.6\% [60.1, 67.1], Llama-3.1-8B 58.4\% [54.8, 62.0]). Detection at the
injected step is rare in all three (21/9/0 flags), so most switches move with the corruption
rather than catching it. This does not score task success, but it converts deviation into
harm-direction evidence: nearly two-thirds of critical-arm switches adopt the injected value
or change action where the design anticipated harm, while explicit skepticism is the exception.

\begin{table}[pos=htbp]
  \centering
  \caption{Wrong-direction endpoint on the 1{,}964 critical-arm switches. Wrong-direction is
  the exclusive action-switch or silent-propagation class; caught is flag-only. Intervals are
  Wilson 95\%.}
  \label{tab:wrongdir}
  \small
  \begin{tabular}{lrrr}
    \toprule
    \textbf{Sample} & \textbf{Switches} & \textbf{Wrong-direction} & \textbf{Caught} \\
    \midrule
    Pooled       & 1{,}964 & 1{,}258 (64.1\% [61.9, 66.1]) & 30 (1.5\% [1.1, 2.2]) \\
    $s=2$        & 911     & 598 (65.6\% [62.5, 68.7])     & 14 (1.5\%) \\
    $s=3$        & 1{,}053 & 660 (62.7\% [59.7, 65.5])     & 16 (1.5\%) \\
    Qwen3.5-9B     & 516   & 374 (72.5\% [68.5, 76.2])     & 21 (4.1\%) \\
    Ministral-3-8B & 726   & 462 (63.6\% [60.1, 67.1])     & 9 (1.2\%) \\
    Llama-3.1-8B   & 722   & 422 (58.4\% [54.8, 62.0])     & 0 (0.0\%) \\
    \bottomrule
  \end{tabular}
\end{table}

\subsection{Task-utility endpoint (exploratory)}

Deviation from the donor is not task error, so P4.6 added a task-utility stream with no new
generations (GPU 0.0\,h; CPU-only reanalysis plus blinded coding;
\texttt{scripts/16\_task\_utility.py}; exploratory under A-010; design logged in D-044 before coding;
\texttt{artifacts/analysis/new\_evidence.json}). It has two layers. First, deterministic
donor-path adherence on all 1{,}039 semantic $s=3$-switched pairs with parsed decisions at
both arms (primary sample; 14 unparsed pairs auto-coded degraded and excluded from agent
coding): each decision is scored against the donor factual continuation (tool-name match for
the 245 tool-call donors, answer-type match for the 794 free-text donors). Adherence is
flat: 29.1\% [26.4, 31.9] at $s=0$ versus 28.9\% [26.2, 31.7] at $s=3$ (paired exact
McNemar $p=.89$; 53.5\% on tool donors both arms; 21.5\% versus 21.3\% on text donors).
The donor path is a narrow reference both arms usually miss, so the cheap deterministic
proxy is insensitive: the corruption changes \emph{which} non-donor decision is taken, not
the match rate.

Second, blinded adequacy coding of the next decision against the TRUE tool result and the
user request on 120 distinct prefixes (40/model, parsed at both arms) $\times$ 2 arms
(240 codes), plus 40 blinded re-codes. Packets carry a blind id only (model, arm, and donor
hidden); the rubric (D-044) is 2~=~adequate, 1~=~degraded (no incorrect state change),
0~=~harmful (incorrect state change or false task-critical information), 9~=~unjudgeable.
Coding is single-agent, author-overseen (T6 forbids subagents, so no multi-session panel;
weaker than the D-043 panel), with the re-code pass in the same session; the re-code
agreement is therefore intra-rater, not independent replication.

On the 115 judgeable pairs, adequate decisions fall from 55.7\% [46.5, 64.4] at $s=0$ to
28.7\% [21.2, 37.6] at $s=3$, and harmful decisions rise from 15.7\% [10.1, 23.4] to
38.3\% [29.9, 47.4] (paired adequate-vs-not McNemar $b=42$, $c=11$,
$p=2.2\times10^{-5}$; Table~\ref{tab:utility}). The adequate drop is largest for Ministral
(60.5\% to 21.1\%; harmful 18.4\% to 50.0\%) and smallest for Llama (34.2\% to 26.3\%;
harmful 21.1\% to 23.7\%); Qwen falls from 71.8\% to 38.5\% adequate (harmful 7.7\% to
41.0\%). Utility tracks the stage-17 classes: among $s=3$ switches coded wrong-direction,
77.8\% (63/81) [67.6, 85.5] are degraded or harmful at $s=3$, versus 55.9\% (19/34)
[39.5, 71.1] among caught/ambiguous. Donor-match and adequacy barely overlap (8 of 115
pairs both match the donor and code adequate), confirming the deterministic proxy carries
almost no utility signal. Re-coding agrees on 36 of 38 items ($\kappa=.92$; 2 further
pairs excluded for a 9).

Limits. Baseline adequacy is modest (55.7\% adequate and 15.7\% harmful at $s=0$), so the
injected error roughly doubles harm rather than creating it from nothing. The donor path is
one valid path and donor-trajectory success was not verified; text-donor type-match is
coarse. Adequacy is judged on the next decision only, not rollout reward. The single-agent,
same-session protocol cannot support population claims; the Wilson intervals are
measurement uncertainty on this sample, not a general harm rate.

\begin{table}[pos=htbp]
  \centering
  \caption{Task-utility endpoint (exploratory; \texttt{artifacts/analysis/new\_evidence.json}).
  Agent adequacy on 115 judgeable $s=0$/$s=3$ pairs (120 sampled; 5 pairs carry a 9) and
  deterministic donor-path adherence on all 1{,}039 switched pairs. Intervals are Wilson
  95\%; McNemar is the paired exact test.}
  \label{tab:utility}
  \small
  \begin{tabular}{lrrrr}
    \toprule
    \textbf{Sample} & \textbf{n} & \textbf{Adequate} & \textbf{Degraded} & \textbf{Harmful} \\
    \midrule
    Agent $s=0$ & 115 & 55.7\% [46.5, 64.4] & 28.7\% [21.2, 37.5] & 15.7\% [10.1, 23.4] \\
    Agent $s=3$ & 115 & 28.7\% [21.2, 37.6] & 33.0\% [25.1, 42.1] & 38.3\% [29.9, 47.4] \\
    \midrule
    \textbf{Sample} & \textbf{n} & \textbf{Match $s=0$} & \textbf{Match $s=3$} & \textbf{McNemar $p$} \\
    \midrule
    Donor adherence & 1{,}039 & 29.1\% [26.4, 31.9] & 28.9\% [26.2, 31.7] & .89 \\
    \bottomrule
  \end{tabular}
\end{table}

\paragraph{Recovery}.\

Recovery is reported as crude cumulative incidence Events/n (primary), with
Kaplan--Meier $1-S_{\mathrm{KM}}(4)$ as a secondary, RMST over the four-step horizon, and an
exponential half-life only if G4 passes (Tables~\ref{tab:recovery} and~\ref{tab:recovery_all},
Figure~\ref{fig:recovery}). The v2 re-estimation rebuilds the risk set from the semantic
switch definition (\texttt{signatures\_differ}, parser v4 / signature v2) instead of the
run-time hash inequality: a prefix whose injected decision differs from the verbatim-V
reference and that did not flag at injection. The run-time (hash) risk set admitted
final-answer paraphrases the semantic coding treats as unchanged---36 Llama, 53 Ministral,
and 43 Qwen transient-$s3$ members (overlap table in
\texttt{artifacts/analysis/g4\_recovery\_v2.md})---and those members recovered at
41.9--88.9\%, inflating the reported curve. Dropping them leaves semantic risk sets of
376/353/221 transient-$s3$ prefixes (Llama/Ministral/Qwen) against run-set sizes 412/406/264
and reported switch counts 377/390/286; the hash-based crude triple 72/264, 141/406, 112/412
in Table~\ref{tab:gate-status} decomposes exactly into semantic plus hash-only events
(54+18, 116+25, 80+32). Censoring is at the last observed step (A-009). G4 does not
pass on the primary set: 24.4\%, 32.9\%, and 21.3\% of the semantic transient $s=3$ risk set
recover by $H=4$ (Qwen, Ministral, Llama; Wilson 95\% intervals [19.2, 30.5], [28.2, 37.9],
[17.4, 25.7]). The KM secondary reaches 28.2/44.6/22.0\%; the gap to crude is the
competing-risk bias from censoring terminal non-recovery states (final answers, parse
failures) as non-informative, and is largest for Ministral, where 118/353 sequences end in a
parse failure and 58/353 in a final answer. The recovery parse-failure rate is 5.3/33.4/5.0\%
(Llama/Ministral/Qwen transient $s3$); only 61/353 Ministral sequences run to $H=4$ without an
event or an early terminal state. Treating final answers and parse failures as non-recoveries
is the crude itself (co-primary); counting the 0--6 semantic returns missed by hash-exact R1
coding moves crude by at most 1.7 points, and the never-run bounds (33 Ministral and 48 Qwen
semantic switchers predate the rescored parser and have no sequences) stay below 39\%. The only
bound that crosses 50\% is the extreme parse-as-recovery counterfactual for Ministral (66.3\% on
the primary transient-$s3$ cell), which would require every generation failure to be a recovery
and is reported as an uninformative upper bound, not an estimate. The
secondary sets reach crude 25.4/33.1/22.9\% (transient $s=2$) and 31.7/35.3/23.9\% (persistent
$s=3$). The persistent arm is
near-inert (the re-applied false value recurred in only a handful of steps, and the prompts are
byte-identical to the transient arm for 97--100\% of shared steps), so those groups are
descriptive and not a persistence contrast. The Weibull shape rejects the exponential assumption
for all three primary groups ($p=.0047$, $<10^{-4}$, $.0006$), and the fitted half-lives
(4.7--10.9 steps across groups) exceed the horizon in
every group. Endpoint composition differs by model: 78/80 Llama recoveries are R1 (signature
reversion), 43/54 Qwen, but only 38/116 Ministral---most Ministral recoveries are explicit flags
(78), not reversion. The marginal log-rank tests are descriptive, not within-prefix inference:
the risk sets share prefixes, so independence fails by design. They find no difference between
the transient and persistent curves (pooled $p=.49$; model-stratified $p=.47$) or between $s2$
and $s3$ ($p=.52$; stratified $p=.52$). The within-prefix check is the exact paired McNemar on
the $H=4$ recovery indicator: $s2$-vs-$s3$ per-model $p=.33/.46/.85$ (Llama/Ministral/Qwen),
transient-vs-persistent $p=1.0/.625/1.0$. Neither the marginal nor the paired test rejects
equality in either contrast
(\texttt{artifacts/analysis/inference\_robustness.json}, recomputed from
\texttt{artifacts/analysis/recovery\_deep.json}). The gate fails on every point estimate of the primary set, and the ``Ministral close'' reading of the KM secondary is retired: Ministral's gate statistic is crude 32.9\% [28.2, 37.9], not KM 44.6\%. The curve and RMST stand; the
half-life is labelled exploratory. The curves flatten before $H=4$; recovery estimates trace to
\texttt{artifacts/analysis/g4\_recovery\_v2.json}.

\begin{table}[pos=htbp]
  \centering
  \caption{Recovery on the semantic risk sets at $H=4$ (v2 re-estimation;
  \texttt{artifacts/analysis/g4\_recovery\_v2.json}). Recovered is crude Events/n (primary);
  KM is the Kaplan--Meier secondary.}
  \label{tab:recovery}
  \footnotesize
  \setlength{\tabcolsep}{4pt}
  \begin{tabular}{llrrrrrrr}
    \toprule
    \textbf{Model} & \textbf{Risk set} & $n$ & \textbf{Events} & \textbf{Censored} & \textbf{Recovered (crude)} & \textbf{KM} & \textbf{RMST} & \textbf{$t_{1/2}$ (Weibull $p$)} \\
    \midrule
    Qwen3.5-9B     & $s=2$ & 185 & 47  & 138 & 25.4\% & 28.9\% & 3.11 & 7.7 (.019) \\
                   & $s=3$ & 221 & 54  & 167 & 24.4\% & 28.2\% & 3.16 & 8.0 (.0047) \\
    Ministral-3-8B & $s=2$ & 311 & 103 & 208 & 33.1\% & 47.7\% & 2.82 & 4.7 ($<10^{-4}$) \\
                   & $s=3$ & 353 & 116 & 237 & 32.9\% & 44.6\% & 2.87 & 5.0 ($<10^{-4}$) \\
    Llama-3.1-8B   & $s=2$ & 345 & 79  & 266 & 22.9\% & 23.6\% & 3.31 & 10.0 (.0008) \\
                   & $s=3$ & 376 & 80  & 296 & 21.3\% & 22.0\% & 3.37 & 10.9 (.0006) \\
    \bottomrule
  \end{tabular}
\end{table}

\begin{table}[pos=htbp]
  \centering
  \caption{Recovery by $H=4$ across model and risk-set groups (Kaplan--Meier secondary;
  crude primary in Table~\ref{tab:recovery}; semantic risk sets).}
  \label{tab:recovery_all}
  \small
  \begin{tabular}{lrrr}
    \toprule
    \textbf{Model} & \textbf{Transient $s=2$} & \textbf{Transient $s=3$} & \textbf{Persistent $s=3$} \\
    \midrule
    Qwen3.5-9B     & 28.9\% (185) & 28.2\% (221) & 35.2\% (63) \\
    Ministral-3-8B & 47.7\% (311) & 44.6\% (353) & 49.3\% (85) \\
    Llama-3.1-8B   & 23.6\% (345) & 22.0\% (376) & 24.9\% (92) \\
    \bottomrule
  \end{tabular}
\end{table}

\begin{figure}[pos=htbp]
  \centering
  \figfile[0.98\linewidth]{figures/fig9.png}
  \caption{Kaplan--Meier survival (secondary; y-axis from 0.55) on the transient semantic
  risk sets by model (Table~\ref{tab:recovery}); curves are the v2 re-estimation in
  \texttt{artifacts/analysis/g4\_recovery\_v2.json}.}
  \label{fig:recovery}
\end{figure}

\paragraph{Sensitivity and heterogeneity}.\

The pre-registered sensitivity analyses exclude the placebo and fit $s=1,\dots,3$; stratify by
domain and donor; vary strict versus lenient action coding; and repeat a 10\% seed sample for
stability (Table~\ref{tab:sensitivity}, Figure~\ref{fig:sensitivity}). The pre-registered sensitivity analyses are complete. The seed-repeat check: 80
prefixes $\times$ 5 arms $\times$ 3 models were regenerated with a different sampling seed and
compared with the base matrix, giving 100.0\% (Llama-3.1-8B), 98.95\% (Ministral-3-8B), and 98.94\%
(Qwen3.5-9B) signature stability, all above the 0.90 gate. The denominators are 385/381/376 of 400
repeated generations; unparsed generations are excluded. The placebo-excluded slope is 1.03
[0.88, 1.19]; task-clustered resampling leaves the pooled interval essentially unchanged
[0.758, 0.996]; the held-out reserve gives $\hat\beta = 0.781$ [0.443, 1.194]. Domain and donor
strata are reported with the diagnostics. Heterogeneity is reported as an interval, not as a
ranking, and no model is declared more or less sensitive than another. Sensitivity fits trace
to the same artifact, with every point estimate positive.

\begin{table}[pos=htbp]
  \centering
  \caption{Pooled slope under pre-registered sensitivity codings.}
  \label{tab:sensitivity}
  \small
  \begin{tabular}{lcc}
    \toprule
    \textbf{Analysis} & $\hat\beta$ & \textbf{95\% CI} \\
    \midrule
    Primary ($s=0,\dots,3$) & 0.877 & [0.770, 0.998] \\
    Strict signature        & 0.876 & [0.768, 0.996] \\
    Lenient action only     & 0.434 & [0.296, 0.574] \\
    Exclude placebo ($s=1,\dots,3$) & 1.031 & [0.881, 1.193] \\
    \bottomrule
  \end{tabular}
\end{table}

\begin{figure}[pos=htbp]
  \centering
  \figfile[0.98\linewidth]{figures/fig10.png}
  \caption{Pooled slope under sensitivity codings and within strata.}
  \label{fig:sensitivity}
\end{figure}



\paragraph{Conditional interpretation}.\

The study is designed to answer one question at one level of resolution: within a frozen real
prefix, does changing the content of a tool observation change the next decision, and is the
change ordered by severity? If the pooled slope is positive and at least two of three models
replicate the sign, the result is a measured shift among movers for three open-weight
8--9B families on $\tau$-bench airline and retail, under one renderer and one severity rubric. It
is not a universal constant, and because three gate criteria failed, no causal or content-sensitivity
claim is made. If the slope is flat, the result is a measurement: at the
achieved precision, plausible content errors of this kind were not associated with a different
next decision in this
sample. Both outcomes are informative, and the second is the one that practitioners are most
likely to misread, because a flat slope bounds how much observation-level checking is worth
relative to tool selection, planning, or policy adherence.

Two statistical properties bound the interpretation. First, $\beta$ is a conditional log-odds,
not a risk difference, and it is not collapsible: the marginal paired contrasts (+10.2, +16.3,
and +6.2 percentage points for $s2-s0$, $s3-s0$, and $s3-s2$) answer the population-average
question and are reported separately from the within-stratum slope. Second, the assigned severity
is measured rather than perfectly manipulated. The audit finds 7--16 of 60 donor-unused items in the critical
arms and 3/30 decision-changing items in $s1$, so $\beta$ is a weighted contrast between intended and
realized doses, not a clean per-level effect parameter; the bounded sensitivity in
\texttt{artifacts/analysis/severity\_sensitivity.json} shows correction steepens the contrast. The sign is stable across every audit pass
and coding, but the per-level magnitudes should not be read as a calibrated dose--response law. The decomposition changes the interpretation of a positive slope. A slope driven by detection is
a slope in which the agent notices the conflict; a slope driven by silent propagation is a slope
in which the agent absorbs the error. The two have opposite engineering implications even though
they produce the same value of $Y$. That is why the paper reports $Y^{U}$, $Y^{F}$, and $Y^{A}$
separately and why the strict silent-propagation form excludes cases that also re-query or flag.

Recovery adds a temporal dimension that a single-step switch rate cannot supply. If most of the
risk set recovers within four steps, a single verification step may be sufficient in practice. If
recovery is slow or censored, the error persists across the episode and the failure is closer to
a state corruption than to a momentary lapse. The four-step horizon is short, and the paper will
say so; a censored curve is a lower bound on persistence, not a claim that the agent never
recovers.

\paragraph{Engineering implications}.\

The result will inform choices at the observation boundary. If wrong content is associated with different decisions and
silent propagation is common, builders should add verification for critical fields, require a
second tool call before acting on a suspicious value, surface a confidence flag, or route
high-stakes actions to human approval. If the slope is flat, those checks may be less urgent than
tool selection, planning, or policy adherence. The study does not recommend a specific product
change; it estimates the size of the problem and the conditions under which it appears. The
transient/persistent contrast was the design's route to the persistence question, but the
persistent re-application was near-inert in execution: the corrupted field rarely recurred in
later observations, so the contrast is not identified here and is left to future work with a
state-aware re-injection.

\paragraph{Evaluation practice}.\

End-to-end success is a useful but coarse metric. Two agents with the same task score can differ
sharply in whether they notice a wrong observation. An observation-level sensitivity metric
complements task success and can be reported alongside it. ToolResult-800 provides a reproducible
test bed for that metric: the same frozen prefixes, arms, renderer, and coding rules can be
applied to new models without re-harvesting trajectories, which makes comparisons across model
families more interpretable than a single aggregate score. The dataset is also a null-result
asset: a well-powered flat slope is evidence about where reliability effort does not pay off.

\paragraph{What would change the claim}.\

Four outcomes would change the claim rather than weaken it. A pooled slope at or below 0.1 logit
per level, or a negative slope, converts the paper to a measurement result. Placebo jitter above
20\% means the $s=0$ arm is not stable and the severity contrast is not identified. An $s=3$
switch rate above 80\% means the top level is a ceiling and severity is not graded; below 10\%
means the manipulation is too weak in this population. Opposite signs across models, or a
positive slope in fewer than two of three, fail the replication gate. Each of these is
pre-registered, so none of them can be reinterpreted after the fact as a discovery.



The limitations are known in advance and are part of the claim boundary. $\tau$-bench tasks are
public and may be memorized by target models, so baseline reliance could be inflated; the
within-prefix contrast remains valid, but the absolute switch rate may not transfer. The donor
prefixes come from gpt-4o and claude-3.5-sonnet, so the target models continue off-policy
histories; the design fixes prefixes across arms and stratifies by donor, but it does not make the
history on-policy. Severity is a constructed ordinal variable. If $s=1$ is associated with as many switches as
$s=2$, the dose interpretation fails, and the paper reports that finding instead of relabelling
the levels.

The study covers three 8--9B model families, two domains, one renderer, deterministic tools, and a
four-step horizon. Only one deterministic prose renderer is used (renderer \texttt{toolresult-prose-v1}),
so the severity slope is template-local: within one fixed template, the treatment bundles content,
format, length, and salience, and the 20.3--45.6\% G1 format effect (main-run $s=0$ rates
19.8/31.9/38.1\%) is the visible part of that bundle. A second renderer on a 100--200 prefix
subsample would test whether the slope is template-specific and is left to future work, as is the
held-out human audit. It does not claim end-to-end task success, frontier-model behaviour, live-API
validity, or a universal law. The tools are benchmark implementations against a pinned database,
not live services, so a wrong observation here is a controlled counterfactual rather than a
production incident. The annotation was produced by blind rater sessions operated with an AI
agent; the reported round is label-free (A-008) and was checked by a multi-role agent panel,
author-overseen, not by
human coauthors or an external vendor. The G0 40-item harvest audit is likewise an AI pre-audit
(40/40) adopted by the author with human verification deferred (D-021, A-005); the freeze stands
provisional on that item. That limits the independence of the audit and is stated in
the declarations. Agreement statistics measure consistency among sessions, not correctness: the
recovery panel found 14/80 rater over-credits while reproducing the automatic coding on 80/80
sequences, and two severity sessions were identical on 116/120 items. The 18~GB slice runs at
batch 4 rather than the planned batch 8; the G2 cap was not approached (about 8.9 GPU-hours of the
45-hour budget), so the sample is complete. Three gate criteria failed and are reported as
failures: the protocol parser-failure rate is 1.55\% against the 1\% threshold (3.45\% if refusals are
counted; A-011), the severity audit's $s2/s3$ decision-changing share is 88.3\% against 90\%, and G4
recovery is crude 27.3/34.7/27.2\% (Kaplan--Meier 22.0--44.6\% secondary) against the 50\% threshold. The
automatic labeling was
corrected once, after the blind audit found a defect (D-028, A-006); the pre-fix analysis is also
reported
above. The G1 placebo criterion was amended post-freeze under A-004 (not blind to the G1 values,
D-041); the analysis gates were not re-thresholded after the confirmatory run.

\section{Conclusion}\label{sec:conclusion}

This paper asks whether the content of a real tool observation is associated with a different
next decision by an agent, whether that association is ordered by the severity of the content
error, and how quickly the agent returns to the reference decision. The design holds 800 real
$\tau$-bench prefixes fixed, varies only the observation across five arms, and reports a
within-prefix conditional-logistic slope with a silent-propagation decomposition and a four-step
recovery curve. The study is pre-registered, gated, and ran at \$0
cash on a university H200 slice. The base matrix, recovery sequences, and seed-repeat shard are
complete; the pooled slope is $0.877$ [0.770, 0.998], positive in all three models, and at least 21.3\% of
critical-arm switches silently repeat the injected value (recall-adjusted $\approx$28.5\%). Three gate criteria failed: the
protocol parser-failure rate
(1.55\% against 1\%), the severity audit's $s2/s3$ decision-changing share (88.3\% against 90\%),
and recovery (crude 27.3/34.7/27.2\%, Kaplan--Meier 22.0--44.6\%, against 50\%); all are reported
as failures. The
contribution is the dataset, the within-prefix design, the audit procedure that caught a labeling
defect and corrected it before the reported analysis, and a measured shift among movers with its
stated limits; because three gate criteria failed, no causal or content-sensitivity claim is made.

\section*{Declarations}
\begin{itemize}
  \setlength{\itemsep}{0pt}
  \setlength{\parsep}{0pt}
  \item \textbf{Funding:} This research received no external funding. The university provides the
  H200 MIG 1g.18gb slice in kind; the study has no cash budget and no paid API.
  \item \textbf{Conflicts of Interest:} The authors declare that they have no known competing
  financial interests or personal relationships that could have appeared to influence the work
  reported in this paper.
  \item \textbf{Data availability:} Raw $\tau$-bench trajectories are not redistributed. The
  upstream benchmark is MIT licensed and is pinned by commit and file hash in
  \texttt{data/manifests/source\_pin.json}. ToolResult-800 releases only scrubbed derived
  data under CC BY 4.0 with upstream notices preserved: frozen prefix payloads, arm observations,
  perturbation records, renderer version, scrubbed ledger extracts, audit packets and sheets,
  analysis JSONs, and hashes; the v1.0 package layout, table-to-artifact map, and one-command
  validation are in \texttt{release/ToolResult-800-v1.0} (\texttt{TABLE\_MAP.md},
  \texttt{VALIDATION.md}). Run logs, weights, caches, and credentials are never released.
  The frozen manifest and the release checklist are the authoritative record. The DOI is pending
  (deposit staged; \texttt{docs/submission/doi\_record.md}) and will be recorded here and in
  \texttt{CITATION.cff} at release.
  \item \textbf{Code availability:} The harvest, replay, rendering, inference, coding, and
  analysis code is available at \url{https://github.com/Atik203/tool-observation} under the MIT
  license. Documentation and derived data are CC BY 4.0.
  \item \textbf{Ethical approval:} The study collects no human-subject data. Annotation was
  produced by blind rater sessions operated with an AI agent and checked by a multi-role agent
  panel, author-overseen; no human annotation was used and direct identifiers are scrubbed
  before freeze. The institutional ethics determination for the annotation procedure is pending
  and will be cited here once issued (tracked in \texttt{docs/submission/checklist.md}).
  \item \textbf{Author contributions:} Conceptualization, data curation, formal analysis,
  investigation, methodology, project administration, software, validation, visualization, and
  writing (original draft) by M.A.R.; supervision, validation, and writing (review and editing)
  by O.T. The
  CRediT allocation was reconciled to the byline and must be confirmed by every
  author before submission.
  \item \textbf{Use of AI-assisted tooling:} Blind rating sessions were operated with an AI agent
  (three sessions for the severity audit, two each for the outcome and recovery audits); a multi-role
  agent panel, author-overseen, re-derived every rating (one verifier per audit), and a second-rater pair
  re-rated a stratified 120-item subset. An AI-assisted mechanical pre-audit was used for the
  40-item G0 check under D-021. The first, label-visible round was accepted by the author on
  2026-11-18 (D-035); the reported label-free round is primary (A-008) and its panel record is
  documented under D-043 (\texttt{artifacts/audit/verification/panel\_record.json}). Rater sessions ran in this
  project's agent-harness interface and wrote the released packet/sheet CSVs; the rater
  instruction files state the audit gates but not the study hypotheses, and the label-free packets
  exclude arm labels, reference decisions, and automatic answers, so raters were blind to the
  hypotheses tested. All sessions share one underlying model family (the session model identity was
  not logged), so their agreement is reproducibility evidence, not human
  independence. Dates: severity/outcome/recovery packets 2026-11-19, panel record 2026-11-18/19,
  task-utility packets and codes 2026-11-22. The task-utility endpoint is exploratory (A-010),
  single-agent coded, same-session intra-rater only ($\kappa=.92$ on 38 re-codes), author-overseen (T6 forbids
  subagents, so no panel was possible there), and its packets, sealed maps, and codes are released
  in \texttt{packets/}. The authors remain responsible for every label, every gate decision, and every
  claim in this paper.
\end{itemize} \section*{Acknowledgements}
The authors thank the curators of $\tau$-bench for releasing the trajectories, task databases, and
harness that make a replayed fixed-prefix design possible, and the university high-performance
computing staff for the H200 MIG allocation. \hbadness=10000
\sloppy
\setlength{\bibsep}{0pt}
\renewcommand{\bibfont}{\fontsize{8.5pt}{9pt}\selectfont}
\bibliographystyle{elsarticle-num}
\bibliography{ref} \vspace{4pt}
\end{document}
